% Kommentare für den Editor (TexWorks/TexMakerX)
% !TeX encoding   = utf8
% !TeX spellcheck = de-DE

% --- LaTeX Vorlage ------------------------------------------------------
% speziell für einfache Dokumente wie Praktikumsprotokolle
%
% Autor: Matthias Pospiech (matthias@pospiech.eu)
% ------------------------------------------------------------------------

% Dokumentenklasse (Koma Script) -----------------------------------------
\documentclass[%
   %draft,     % Entwurfsstadium
   final,      % fertiges Dokument
   paper=a4, paper=portrait, pagesize=auto, % Papier Einstellungen
   fontsize=11pt, % Schriftgröße
   ngerman, % Sprache 
 ]{scrreprt} % Classes: scrartcl, scrreprt, scrbook

% ~~~~~~~~~~~~~~~~~~~~~~~~~~~~~~~~~~~~~~~~~~~~~~~~~~~~~~~~~~~~~~~~~~~~~~~~
% encoding
% ~~~~~~~~~~~~~~~~~~~~~~~~~~~~~~~~~~~~~~~~~~~~~~~~~~~~~~~~~~~~~~~~~~~~~~~~

% Encoding der Dateien (sonst funktionieren Umlaute nicht)
\usepackage[utf8]{inputenc}

% Encoding der Verzeichnisse (für Pfade mit Umlauten und Leerzeichne)
\usepackage[%
   extendedchars, encoding, multidot, space,
   %filenameencoding=latin1, % Windows XP, Vista, 7
   filenameencoding=utf8,   % Linux, OS X
]{grffile}

% ~~~~~~~~~~~~~~~~~~~~~~~~~~~~~~~~~~~~~~~~~~~~~~~~~~~~~~~~~~~~~~~~~~~~~~~~
% Pakete und Stile
% ~~~~~~~~~~~~~~~~~~~~~~~~~~~~~~~~~~~~~~~~~~~~~~~~~~~~~~~~~~~~~~~~~~~~~~~~
% Schriften
% ~~~~~~~~~~~~~~~~~~~~~~~~~~~~~~~~~~~~~~~~~~~~~~~~~~~~~~~~~~~~~~~~~~~~~~~~
% Fonts Fonts Fonts
% ~~~~~~~~~~~~~~~~~~~~~~~~~~~~~~~~~~~~~~~~~~~~~~~~~~~~~~~~~~~~~~~~~~~~~~~~

% immer laden:
\usepackage[T1]{fontenc} % T1 Schrift Encoding
\usepackage{textcomp}	 % Zusätzliche Symbole (Text Companion font extension)

% ~~~~~~~~~~~~~~~~~~~~~~~~~~~~~~~~~~~~~~~~~~~~~~~~~~~~~~~~~~~~~~~~~~~~~~~~
% Symbole
% ~~~~~~~~~~~~~~~~~~~~~~~~~~~~~~~~~~~~~~~~~~~~~~~~~~~~~~~~~~~~~~~~~~~~~~~~

\usepackage{amssymb}
\usepackage{mathcomp}


%% ==== Zusammengesetzte Schriften  (Sans + Serif) =======================

%% - Latin Modern
\usepackage{lmodern}
%% -------------------

%% - Bera Schriften
%\usepackage{bera}
%% -------------------

%% - Times, Helvetica, Courier (Word Standard...)
%\usepackage{mathptmx}
%\usepackage[scaled=.90]{helvet}
%\usepackage{courier}
%% -------------------

%% - Palantino , Helvetica, Courier
%\usepackage{mathpazo}
%\usepackage[scaled=.95]{helvet}
%\usepackage{courier}
%% -------------------

%% - Charter, Bera Sans
%\usepackage{charter}\linespread{1.05}
%\renewcommand{\sfdefault}{fvs}
%\usepackage[charter]{mathdesign}



%%%% =========== Typewriter =============

%\usepackage{courier}                   %% --- Courier
%\renewcommand{\ttdefault}{cmtl}        %% --- CmBright Typewriter Font
%\usepackage[%                          %% --- Luxi Mono (Typewriter)
%   scaled=0.9
%]{luximono}



% Pakete Laden
% ~~~~~~~~~~~~~~~~~~~~~~~~~~~~~~~~~~~~~~~~~~~~~~~~~~~~~~~~~~~~~~~~~~~~~~~~
% These packages must be loaded before all others
% (primarily because they are required by other packages)
% ~~~~~~~~~~~~~~~~~~~~~~~~~~~~~~~~~~~~~~~~~~~~~~~~~~~~~~~~~~~~~~~~~~~~~~~~
\usepackage{calc}
\usepackage{fixltx2e}	% Fix known LaTeX2e bugs

\usepackage[ngerman]{babel} 	% Sprache
\usepackage[dvipsnames, table]{xcolor} 	% Farben
\usepackage[ngerman, iso]{isodate} 	%Datumsformat
\usepackage[utf8]{inputenc}

% ~~~~~~~~~~~~~~~~~~~~~~~~~~~~~~~~~~~~~~~~~~~~~~~~~~~~~~~~~~~~~~~~~~~~~~~~
% Bilder, Gleitumgebungen und Platzierung
% ~~~~~~~~~~~~~~~~~~~~~~~~~~~~~~~~~~~~~~~~~~~~~~~~~~~~~~~~~~~~~~~~~~~~~~~~

\usepackage[]{graphicx}					% Graphiken
\usepackage{epstopdf}		% konvertiert eps in pdf

% provides new floats and enables H float modifier option
\usepackage{float}
% Floats immer erst nach der Referenz setzen
\usepackage{flafter}
% Alel Floats werden vor der nächsten section ausgegeben
\usepackage[section]{placeins} 
%

\usepackage{tikz}
%tikz verwenden um Grafiken zu erstellen
\usepackage{pgfplots}
%pgfplots verwenden für Diagramme

\usepackage[european, siunitx]{circuitikz}
% Schaltpläne

% ~~~~~~~~~~~~~~~~~~~~~~~~~~~~~~~~~~~~~~~~~~~~~~~~~~~~~~~~~~~~~~~~~~~~~~~~
% Beschriftungen (captions)
% ~~~~~~~~~~~~~~~~~~~~~~~~~~~~~~~~~~~~~~~~~~~~~~~~~~~~~~~~~~~~~~~~~~~~~~~~

\usepackage{caption}
\usepackage{subcaption}

% ~~~~~~~~~~~~~~~~~~~~~~~~~~~~~~~~~~~~~~~~~~~~~~~~~~~~~~~~~~~~~~~~~~~~~~~~
% Math
% ~~~~~~~~~~~~~~~~~~~~~~~~~~~~~~~~~~~~~~~~~~~~~~~~~~~~~~~~~~~~~~~~~~~~~~~~

% Base Math Package
\usepackage[fleqn]{amsmath} 
% Warnt bei Benutzung von Befehlen die mit amsmath inkompatibel sind.
\usepackage[all, error]{onlyamsmath}
%Brüche schräg statt horizontal brechen
\usepackage{nicefrac}
% mathtools für rcases (geschweifte Klammern rechts)
\usepackage{mathtools}

% trfsigns für laplace
\usepackage{trfsigns}

% ~~~~~~~~~~~~~~~~~~~~~~~~~~~~~~~~~~~~~~~~~~~~~~~~~~~~~~~~~~~~~~~~~~~~~~~~
% Science
% ~~~~~~~~~~~~~~~~~~~~~~~~~~~~~~~~~~~~~~~~~~~~~~~~~~~~~~~~~~~~~~~~~~~~~~~~

% Einheiten und Zahlenformatierung
\usepackage{siunitx}

% ~~~~~~~~~~~~~~~~~~~~~~~~~~~~~~~~~~~~~~~~~~~~~~~~~~~~~~~~~~~~~~~~~~~~~~~~
% Tables (Tabular)
% ~~~~~~~~~~~~~~~~~~~~~~~~~~~~~~~~~~~~~~~~~~~~~~~~~~~~~~~~~~~~~~~~~~~~~~~~

\usepackage{booktabs}
\usepackage{ltxtable} % Longtable + tabularx

% ~~~~~~~~~~~~~~~~~~~~~~~~~~~~~~~~~~~~~~~~~~~~~~~~~~~~~~~~~~~~~~~~~~~~~~~~
% text related packages
% ~~~~~~~~~~~~~~~~~~~~~~~~~~~~~~~~~~~~~~~~~~~~~~~~~~~~~~~~~~~~~~~~~~~~~~~~

\usepackage{url}            % Befehl \url{...}
\usepackage{enumitem}		% Kompakte Listen

% Neue Befehle: \Centering, \RaggedLeft, and \RaggedRight, ... 
\usepackage{ragged2e}


% ~~~~~~~~~~~~~~~~~~~~~~~~~~~~~~~~~~~~~~~~~~~~~~~~~~~~~~~~~~~~~~~~~~~~~~~~
% Citations
% ~~~~~~~~~~~~~~~~~~~~~~~~~~~~~~~~~~~~~~~~~~~~~~~~~~~~~~~~~~~~~~~~~~~~~~~~

%\usepackage[
%	style=alphabetic, % Loads the bibliography and the citation style 
%	natbib=true, % define natbib compatible cite commands
%]{biblatex}	
% Other options:
%	style=numeric, % 
%	style=numeric-comp,    % [1–3, 7, 8]
%	style=numeric-verb,    % [2]; [5]; [6]

\usepackage{csquotes}

% ~~~~~~~~~~~~~~~~~~~~~~~~~~~~~~~~~~~~~~~~~~~~~~~~~~~~~~~~~~~~~~~~~~~~~~~~
% layout packages
% ~~~~~~~~~~~~~~~~~~~~~~~~~~~~~~~~~~~~~~~~~~~~~~~~~~~~~~~~~~~~~~~~~~~~~~~~
%
% Befehle für 1,5 und 2 zeilig: 
% \singlespacing, \onehalfspacing und \doublespacing
\usepackage{setspace}

% ~~~~~~~~~~~~~~~~~~~~~~~~~~~~~~~~~~~~~~~~~~~~~~~~~~~~~~~~~~~~~~~~~~~~~~~~
% Kopf und Fusszeile
% ~~~~~~~~~~~~~~~~~~~~~~~~~~~~~~~~~~~~~~~~~~~~~~~~~~~~~~~~~~~~~~~~~~~~~~~~

% Kopf und Fusszeile mit scrpage2 einstellen
\usepackage[automark, komastyle, nouppercase]{scrpage2}

% ~~~~~~~~~~~~~~~~~~~~~~~~~~~~~~~~~~~~~~~~~~~~~~~~~~~~~~~~~~~~~~~~~~~~~~~~
% pdf packages
% ~~~~~~~~~~~~~~~~~~~~~~~~~~~~~~~~~~~~~~~~~~~~~~~~~~~~~~~~~~~~~~~~~~~~~~~~

% Include pages from external PDF documents in LaTeX documents
\usepackage{pdfpages} 

% Optischer Randausgleich mit pdfTeX
\usepackage{microtype}

%%% svg to tex
\newcommand{\executeiffilenewer}[3]{%
\ifnum\pdfstrcmp{\pdffilemoddate{#1}}%
{\pdffilemoddate{#2}}>0%
{\immediate\write18{#3}}\fi%
}

% Einstellungen und Layoutstile 
% ~~~~~~~~~~~~~~~~~~~~~~~~~~~~~~~~~~~~~~~~~~~~~~~~~~~~~~~~~~~~~~~~~~~~~~~~
% Colors
% ~~~~~~~~~~~~~~~~~~~~~~~~~~~~~~~~~~~~~~~~~~~~~~~~~~~~~~~~~~~~~~~~~~~~~~~~
\definecolor{sectioncolor}{RGB}{0, 0, 0}     % black

% ~~~~~~~~~~~~~~~~~~~~~~~~~~~~~~~~~~~~~~~~~~~~~~~~~~~~~~~~~~~~~~~~~~~~~~~~
% text related 
% ~~~~~~~~~~~~~~~~~~~~~~~~~~~~~~~~~~~~~~~~~~~~~~~~~~~~~~~~~~~~~~~~~~~~~~~~

%% style of URL
\urlstyle{tt}


% Keine hochgestellten Ziffern in der Fussnote (KOMA-Script-spezifisch):
\deffootnote{1.5em}{1em}{\makebox[1.5em][l]{\thefootnotemark}}

% Limit space of footnotes to 10 lines
\setlength{\dimen\footins}{10\baselineskip}

% prevent continuation of footnotes 
% at facing page
\interfootnotelinepenalty=10000 

% ~~~~~~~~~~~~~~~~~~~~~~~~~~~~~~~~~~~~~~~~~~~~~~~~~~~~~~~~~~~~~~~~~~~~~~~~
% Science
% ~~~~~~~~~~~~~~~~~~~~~~~~~~~~~~~~~~~~~~~~~~~~~~~~~~~~~~~~~~~~~~~~~~~~~~~~

\sisetup{%
	mode = math, detect-family, detect-weight,	
	exponent-product = \cdot,
	number-unit-separator=\text{\,},
	output-decimal-marker={,},
}

% ~~~~~~~~~~~~~~~~~~~~~~~~~~~~~~~~~~~~~~~~~~~~~~~~~~~~~~~~~~~~~~~~~~~~~~~~
% Citations / Style of Bibliography
% ~~~~~~~~~~~~~~~~~~~~~~~~~~~~~~~~~~~~~~~~~~~~~~~~~~~~~~~~~~~~~~~~~~~~~~~~

% Kommentar entfernene wenn biblatex geladen wird
% \IfPackageLoaded{biblatex}{%
	\ExecuteBibliographyOptions{%
%--- Backend --- --- ---
	backend=bibtex,  % (bibtex, bibtex8, biber)
	bibwarn=true, %
	bibencoding=ascii, % (ascii, inputenc, <encoding>)
%--- Sorting --- --- ---
	sorting=nty, % Sort by name, title, year.
	% other options: 
	% nty        Sort by name, title, year.
	% nyt        Sort by name, year, title.
	% nyvt       Sort by name, year, volume, title.
	% anyt       Sort by alphabetic label, name, year, title.
	% anyvt      Sort by alphabetic label, name, year, volume, title.
	% ynt        Sort by year, name, title.
	% ydnt       Sort by year (descending), name, title.
	% none       Do not sort at all. All entries are processed in citation order.
	% debug      Sort by entry key. This is intended for debugging only.
	%
	sortcase=true,
	sortlos=los, % (bib, los) The sorting order of the list of shorthands
	sortcites=false, % do/do not sort citations according to bib	
%--- Dates --- --- ---
	date=comp,  % (short, long, terse, comp, iso8601)
%	origdate=
%	eventdate=
%	urldate=
%	alldates=
	datezeros=true, %
	dateabbrev=true, %
%--- General Options --- --- ---
	maxnames=1,
	minnames=1,
%	maxbibnames=99,
%	maxcitenames=1,
%	autocite= % (plain, inline, footnote, superscript) 
	autopunct=true,
	language=auto,
	babel=none, % (none, hyphen, other, other*)
	block=none, % (none, space, par, nbpar, ragged)
	notetype=foot+end, % (foot+end, footonly, endonly)
	hyperref=true, % (true, false, auto)
	backref=true,
	backrefstyle=three, % (none, three, two, two+, three+, all+)
	backrefsetstyle=setonly, %
	indexing=false, % 
	% options:
	% true       Enable indexing globally.
	% false      Disable indexing globally.
	% cite       Enable indexing in citations only.
	% bib        Enable indexing in the bibliography only.
	refsection=none, % (part, chapter, section, subsection)
	refsegment=none, % (none, part, chapter, section, subsection)
	abbreviate=true, % (true, false)
	defernumbers=false, % 
	punctfont=false, % 
	arxiv=abs, % (ps, pdf, format)	
%--- Style Options --- --- ---	
% The following options are provided by the standard styles
	isbn=false,%
	url=false,%
	doi=false,%
	eprint=false,%	
	}%	
	
	% change alpha label to be without +	
	\renewcommand*{\labelalphaothers}{}
	
	% change 'In: <magazine>" to "<magazine>"
	\renewcommand*{\intitlepunct}{}
	\DefineBibliographyStrings{german}{in={}}
	
	% make names capitalized \textsc{}
	\renewcommand{\mkbibnamefirst}{\textsc}
	\renewcommand{\mkbibnamelast}{\textsc}
	
	% make volume and number look like 
	% 'Bd. 33(14): '
	\renewbibmacro*{volume+number+eid}{%
	  \setunit{\addcomma\space}%
	  \bibstring{volume}% 
	  \setunit{\addspace}%
	  \printfield{volume}%
	  \iffieldundef{number}{}{% 
	    \printtext[parens]{%
	      \printfield{number}%
	    }%
	  }%
	  \setunit{\addcomma\space}%
	  \printfield{eid}
	  %\setunit{\addcolon\space}%
	  }	

	% <authors>: <title>
	\renewcommand*{\labelnamepunct}{\addcolon\space}
	% make ': ' before pages
	\renewcommand*{\bibpagespunct}{\addcolon\space}
	% names delimiter ';' instead of ','
	%\renewcommand*{\multinamedelim}{\addsemicolon\space}

	% move date before issue
	\renewbibmacro*{journal+issuetitle}{%
	  \usebibmacro{journal}%
	  \setunit*{\addspace}%
	  \iffieldundef{series}
	    {}
	    {\newunit
	     \printfield{series}%
	     \setunit{\addspace}}%
	  %
	  \usebibmacro{issue+date}%
	  \setunit{\addcolon\space}%
	  \usebibmacro{issue}%
	  \setunit{\addspace}%
	  \usebibmacro{volume+number+eid}%
	  \newunit}

	% print all names, even if maxnames = 1
	\DeclareCiteCommand{\citeauthors}
	  {
	   \defcounter{maxnames}{1000}
	   \boolfalse{citetracker}%
	   \boolfalse{pagetracker}%
	   \usebibmacro{prenote}}
	  {\ifciteindex
	     {\indexnames{labelname}}
	     {}%
	   \printnames{labelname}}
	  {\multicitedelim}
	  {\usebibmacro{postnote}}

}%

% ~~~~~~~~~~~~~~~~~~~~~~~~~~~~~~~~~~~~~~~~~~~~~~~~~~~~~~~~~~~~~~~~~~~~~~~~
% figures, placement, floats and captions
% ~~~~~~~~~~~~~~~~~~~~~~~~~~~~~~~~~~~~~~~~~~~~~~~~~~~~~~~~~~~~~~~~~~~~~~~~

% Make float placement easier
\renewcommand{\floatpagefraction}{.75} % vorher: .5
\renewcommand{\textfraction}{.1}       % vorher: .2
\renewcommand{\topfraction}{.8}        % vorher: .7
\renewcommand{\bottomfraction}{.5}     % vorher: .3
\setcounter{topnumber}{3}        % vorher: 2
\setcounter{bottomnumber}{2}     % vorher: 1
\setcounter{totalnumber}{5}      % vorher: 3

%% ~~~ Captions ~~~~~~~~~~~~~~~~~~~~~~~~~~~~~~~~~~~~~~~~~~~~~~~~~~~~~~~~~~
% Style of captions
\DeclareCaptionStyle{captionStyleTemplateDefault}
[ % single line captions
   justification = centering
]
{ % multiline captions
% -- Formatting
   format      = plain,  % plain, hang
   indention   = 0em,    % indention of text 
   labelformat = default,% default, empty, simple, brace, parens
   labelsep    = colon,  % none, colon, period, space, quad, newline, endash
   textformat  = simple, % simple, period
% -- Justification
   justification = justified, %RaggedRight, justified, centering
   singlelinecheck = true, % false (true=ignore justification setting in single line)
% -- Fonts
   labelfont   = {small,bf},
   textfont    = {small,rm},
% valid values:
% scriptsize, footnotesize, small, normalsize, large, Large
% normalfont, ip, it, sl, sc, md, bf, rm, sf, tt
% singlespacing, onehalfspacing, doublespacing
% normalcolor, color=<...>
%
% -- Margins and further paragraph options
   margin = 10pt, %.1\textwidth,
   % width=.8\linewidth,
% -- Skips
   skip     = 10pt, % vertical space between the caption and the figure
   position = auto, % top, auto, bottom
% -- Lists
   % list=no, % suppress any entry to list of figure 
   listformat = subsimple, % empty, simple, parens, subsimple, subparens
% -- Names & Numbering
   % figurename = Abb. %
   % tablename  = Tab. %
   % listfigurename=
   % listtablename=
   % figurewithin=chapter
   % tablewithin=chapter
%-- hyperref related options
	hypcap=true, % (true, false) 
	% true=all hyperlink anchors are placed at the 
	% beginning of the (floating) environment
	%
	hypcapspace=0.5\baselineskip
}

% apply caption style
\captionsetup{
	style = captionStyleTemplateDefault % base
}

% Predefinded skip setup for different floats
\captionsetup[table]{position=top}
\captionsetup[figure]{position=bottom}


% options for subcaptions
\captionsetup[sub]{ %
	style = captionStyleTemplateDefault, % base
	skip=6pt,
	margin=5pt,
	labelformat = parens,% default, empty, simple, brace
	labelsep    = space,
	list=false,
	hypcap=false
}

% ~~~~~~~~~~~~~~~~~~~~~~~~~~~~~~~~~~~~~~~~~~~~~~~~~~~~~~~~~~~~~~~~~~~~~~~~
% layout 
% ~~~~~~~~~~~~~~~~~~~~~~~~~~~~~~~~~~~~~~~~~~~~~~~~~~~~~~~~~~~~~~~~~~~~~~~~


%% Paragraph Separation =================================
\KOMAoptions{%
   parskip=absolute, % do not change indentation according to fontsize
   parskip=false     % indentation of 1em
   % parskip=half    % parksip of 1/2 line 
}%

%% line spacing =========================================
%\onehalfspacing	% 1,5-facher Abstand
%\doublespacing		% 2-facher Abstand

%% page layout ==========================================

\raggedbottom     % Variable Seitenhoehen zulassen

% Koma Script text area layout
\KOMAoptions{%
   DIV=11,% (Size of Text Body, higher values = greater textbody)
   BCOR=5mm% (Bindekorrektur)
}%

%%% === Page Layout  Options ===
\KOMAoptions{% (most options are for package typearea)
   % twoside=true, % two side layout (alternating margins, standard in books)
   twoside=false, % single side layout 
   %
   headlines=2.1,%
}%

%\KOMAoptions{%
%      headings=noappendixprefix % chapter in appendix as in body text
%      ,headings=nochapterprefix  % no prefix at chapters
%      % ,headings=appendixprefix   % inverse of 'noappendixprefix'
%      % ,headings=chapterprefix    % inverse of 'nochapterprefix'
%      % ,headings=openany   % Chapters start at any side
%      % ,headings=openleft  % Chapters start at left side
%      ,headings=openright % Chapters start at right side      
%}%


% reloading of typearea, necessary if setting of spacing changed
\typearea[current]{last}

%% Zitate
\renewcommand*{\dictumwidth}{0.5\textwidth}

% ~~~~~~~~~~~~~~~~~~~~~~~~~~~~~~~~~~~~~~~~~~~~~~~~~~~~~~~~~~~~~~~~~~~~~~~~
% Titlepage
% ~~~~~~~~~~~~~~~~~~~~~~~~~~~~~~~~~~~~~~~~~~~~~~~~~~~~~~~~~~~~~~~~~~~~~~~~
\KOMAoptions{%
   % titlepage=true %
   titlepage=false %
}%

% ~~~~~~~~~~~~~~~~~~~~~~~~~~~~~~~~~~~~~~~~~~~~~~~~~~~~~~~~~~~~~~~~~~~~~~~~
% head and foot lines
% ~~~~~~~~~~~~~~~~~~~~~~~~~~~~~~~~~~~~~~~~~~~~~~~~~~~~~~~~~~~~~~~~~~~~~~~~

% \pagestyle{scrheadings} % Seite mit Headern
\pagestyle{scrplain} % Seiten ohne Header

% loescht voreingestellte Stile
\clearscrheadings
\clearscrplain
%
% Was steht wo...
% Bei headings:
%   % Oben aussen: Kapitel und Section
%   % Unten aussen: Seitenzahl
%   \ohead{\pagemark}
%   \ihead{\headmark}
%   \ofoot[\pagemark]{} % Außen unten: Seitenzahlen bei plain
% Bei Plain:
\cfoot[\pagemark]{\pagemark} % Mitte unten: Seitenzahlen bei plain


% Angezeigte Abschnitte im Header
% \automark[section]{chapter} %[rechts]{links}
\automark[subsection]{section} %[rechts]{links}

% ~~~~~~~~~~~~~~~~~~~~~~~~~~~~~~~~~~~~~~~~~~~~~~~~~~~~~~~~~~~~~~~~~~~~~~~~
% headings / page opening
% ~~~~~~~~~~~~~~~~~~~~~~~~~~~~~~~~~~~~~~~~~~~~~~~~~~~~~~~~~~~~~~~~~~~~~~~~
\setcounter{secnumdepth}{3}

\KOMAoptions{%
%%%% headings
   % headings=small  % Small Font Size, thin spacing above and below
   % headings=normal % Medium Font Size, medium spacing above and below
   headings=big % Big Font Size, large spacing above and below
}%

% Titelzeile linksbuendig, haengend
\renewcommand*{\raggedsection}{\raggedright} 

% ~~~~~~~~~~~~~~~~~~~~~~~~~~~~~~~~~~~~~~~~~~~~~~~~~~~~~~~~~~~~~~~~~~~~~~~~
% fonts of headings
% ~~~~~~~~~~~~~~~~~~~~~~~~~~~~~~~~~~~~~~~~~~~~~~~~~~~~~~~~~~~~~~~~~~~~~~~~
\setkomafont{sectioning}{\normalfont\sffamily} % \rmfamily
\setkomafont{descriptionlabel}{\itshape}
\setkomafont{pageheadfoot}{\normalfont\normalcolor\small\sffamily}
\setkomafont{pagenumber}{\normalfont\sffamily}

%%% --- Titlepage ---
%\setkomafont{subject}{}
%\setkomafont{subtitle}{}
%\setkomafont{title}{}

% ~~~~~~~~~~~~~~~~~~~~~~~~~~~~~~~~~~~~~~~~~~~~~~~~~~~~~~~~~~~~~~~~~~~~~~~~
% settings and layout of TOC, LOF, 
% ~~~~~~~~~~~~~~~~~~~~~~~~~~~~~~~~~~~~~~~~~~~~~~~~~~~~~~~~~~~~~~~~~~~~~~~~
\setcounter{tocdepth}{3} % Depth of TOC Display

% ~~~~~~~~~~~~~~~~~~~~~~~~~~~~~~~~~~~~~~~~~~~~~~~~~~~~~~~~~~~~~~~~~~~~~~~~
% Tabellen
% ~~~~~~~~~~~~~~~~~~~~~~~~~~~~~~~~~~~~~~~~~~~~~~~~~~~~~~~~~~~~~~~~~~~~~~~~

%%% -| Neue Spaltendefinitionen 'columntypes' |--
%
% Belegte Spaltentypen:
% l - links
% c - zentriert
% r - rechts
% p,m,b  - oben, mittig, unten
% X - tabularx Auto-Spalte

% um Tabellenspalten mit Flattersatz zu setzen, muss \\ vor
% (z.B.) \raggedright geschuetzt werden:
\newcommand{\PreserveBackslash}[1]{\let\temp=\\#1\let\\=\temp}

% Spalten mit Flattersatz und definierte Breite:
% m{} -> mittig
% p{} -> oben
% b{} -> unten
%
% Linksbuendig:
\newcolumntype{v}[1]{>{\PreserveBackslash\RaggedRight\hspace{0pt}}p{#1}}
\newcolumntype{M}[1]{>{\PreserveBackslash\RaggedRight\hspace{0pt}}m{#1}}
% % Rechtsbuendig :
% \newcolumntype{R}[1]{>{\PreserveBackslash\RaggedLeft\hspace{0pt}}m{#1}}
% \newcolumntype{S}[1]{>{\PreserveBackslash\RaggedLeft\hspace{0pt}}p{#1}}
% % Zentriert :
% \newcolumntype{Z}[1]{>{\PreserveBackslash\Centering\hspace{0pt}}m{#1}}
% \newcolumntype{A}[1]{>{\PreserveBackslash\Centering\hspace{0pt}}p{#1}}

\newcolumntype{Y}{>{\PreserveBackslash\RaggedLeft\hspace{0pt}}X}

%-- Einstellungen für Tabellen ----------
\providecommand\tablestyle{%
  \renewcommand{\arraystretch}{1.4} % Groessere Abstaende zwischen Zeilen
  \normalfont\normalsize            %
  \sffamily\small           % Serifenlose und kleine Schrift
  \centering%                       % Tabelle zentrieren
}

%--Einstellungen für Tabellen ----------

\colorlet{tablesubheadcolor}{gray!40}
\colorlet{tableheadcolor}{gray!25}
\colorlet{tableblackheadcolor}{black!60}
\colorlet{tablerowcolor}{gray!15.0}


% ~~~~~~~~~~~~~~~~~~~~~~~~~~~~~~~~~~~~~~~~~~~~~~~~~~~~~~~~~~~~~~~~~~~~~~~~
% pdf packages
% ~~~~~~~~~~~~~~~~~~~~~~~~~~~~~~~~~~~~~~~~~~~~~~~~~~~~~~~~~~~~~~~~~~~~~~~~

% ~~~~~~~~~~~~~~~~~~~~~~~~~~~~~~~~~~~~~~~~~~~~~~~~~~~~~~~~~~~~~~~~~~~~~~~~
% fix remaining problems
% ~~~~~~~~~~~~~~~~~~~~~~~~~~~~~~~~~~~~~~~~~~~~~~~~~~~~~~~~~~~~~~~~~~~~~~~~




% ~~~~~~~~~~~~~~~~~~~~~~~~~~~~~~~~~~~~~~~~~~~~~~~~~~~~~~~~~~~~~~~~~~~~~~~~
% Eigene Befehle
% ~~~~~~~~~~~~~~~~~~~~~~~~~~~~~~~~~~~~~~~~~~~~~~~~~~~~~~~~~~~~~~~~~~~~~~~~
% -- new commands --
\providecommand{\abs}[1]{\lvert#1\rvert}
\providecommand{\Abs}[1]{\left\lvert#1\right\rvert}
\providecommand{\norm}[1]{\left\Vert#1\right\Vert}
\providecommand{\Trace}[1]{\ensuremath{\Tr\{\,#1\,\}}} % Trace /Spur
%

\renewcommand{\d}{\partial\mspace{2mu}} % partial diff
\newcommand{\td}{\,\mathrm{d}}	% total diff

\newcommand{\Ham}{\mathcal{H}}    % Hamilton
\newcommand{\Prob}{\mathscr{P}}    % Hamilton
\newcommand{\unity}{\mathds{1}}   % Real

\renewcommand{\i}{\mathrm{i}}   % imagin�re Einheit



% -- New Operators --
\DeclareMathOperator{\rot}{rot}
\DeclareMathOperator{\grad}{grad}
\DeclareMathOperator{\Tr}{Tr}
\DeclareMathOperator{\const}{const}

%\renewcommand{\labelenumi}{\alph{enumi})}
%\renewcommand{\labelenumii}{\roman{enumii})}
\setlength{\parindent}{0pt}
% ~~~~~~~~~~~~~~~~~~~~~~~~~~~~~~~~~~~~~~~~~~~~~~~~~~~~~~~~~~~~~~~~~~~~~~~~
% Eigene Befehle
% ~~~~~~~~~~~~~~~~~~~~~~~~~~~~~~~~~~~~~~~~~~~~~~~~~~~~~~~~~~~~~~~~~~~~~~~~
% Silbentrennung hinzufügen als 
% Sil-ben-tren-nung 
\hyphenation{}

\listfiles % schreibt alle verwendeten Dateien in die log Datei

%% Dokument Beginn %%%%%%%%%%%%%%%%%%%%%%%%%%%%%%%%%%%%%%%%%%%%%%%%%%%%%%%%
\begin{document}
%\includepdf{content/Deckblatt_WI}
%% Automatische Titelseite

\subject{Mathematik 3}
\title{Formelsammlung\\}
\author{\includegraphics[width=6.5cm]{images/logo}\\ \\  David Mändlen et al.}
\date{\today}
\maketitle

% Manuelle Titelseite

%\begin{titlepage}
%   \mbox{}\vspace{5\baselineskip}\\
%   \sffamily\huge
%   \centering
%   % Titel
%   \LaTeXe{}-Vorlage von Matthias Pospiech
%   \vspace{3\baselineskip}\\
%   \rmfamily\Large
%   Leibniz Universität Hannover
%   \vspace{2\baselineskip}\\
%   \rmfamily\Large
%   Matthias Pospiech
%   \vspace{1\baselineskip}\\
%   \today
%\end{titlepage}


%\clearpage

\null

\vfill

\dictum[Dr. Axel Stoll]{Magie ist Physik durch Wollen!}

\clearpage

%\tableofcontents

% Testdokumente (auskommentieren!)
%\input{content/demo/demo.tex}
%\input{content/demo/latexexample.tex}

% in diese Datei gehört der Inhalt des Dokumentes:
%\input{content/Fehlerrechnung.tex}
%\input{content/Vorbereitung.tex}
% Automatische Titelseite

\subject{Mathematik 3}
\title{Formelsammlung\\}
\author{\includegraphics[width=6.5cm]{images/logo}\\ \\  David Mändlen et al.}
\date{\today}
\maketitle

% Manuelle Titelseite

%\begin{titlepage}
%   \mbox{}\vspace{5\baselineskip}\\
%   \sffamily\huge
%   \centering
%   % Titel
%   \LaTeXe{}-Vorlage von Matthias Pospiech
%   \vspace{3\baselineskip}\\
%   \rmfamily\Large
%   Leibniz Universität Hannover
%   \vspace{2\baselineskip}\\
%   \rmfamily\Large
%   Matthias Pospiech
%   \vspace{1\baselineskip}\\
%   \today
%\end{titlepage}


\tableofcontents
%% Ableitungs und Integral
\chapter{Grundlagen}
\section{Ableitungsregeln}
\begin{table}[h]
\centering
\renewcommand{\arraystretch}{2}
\begin{tabularx}{\textwidth}{llcl}
\textbf{Summenregel} & $f(x) = u(x) + n(x)$ &  $\Rightarrow$ &  $f'(x) = u'(x) + n'(x)$	 \\
\textbf{Produktregel} & $f(x) = u(x) \cdot n(x)$	 &  $\Rightarrow$	 &  $f'(x) = u'(x) \cdot n(x) + u(x) \cdot n'(x)$	 \\
\textbf{Produktregel} & $f(x) = u(x) \cdot v(x) \cdot w(x)$ &  $\Rightarrow$ &  $f'(x) = u'(x) \cdot v(x) \cdot w(x) + $\\
\textbf{für 3 Produkte}	 &  &  &  $u(x) \cdot v'(x) \cdot w(x) + $\\
 &  &  & $u(x) \cdot v(x) \cdot w'(x)$ \\
\textbf{Quotientenregel} & $f(x) = \frac{u(x)}{n(x)}$ &  $\Rightarrow$ &  $f'(x) = \frac{u'(x) \cdot n(x) - u(x) \cdot n'(x)}{(n(x))^2}$ \\
\textbf{Kettenregel} & $f(x) = n(u(x))$ &  $\Rightarrow$	 &  $f'(x) = n'(u(x)) \cdot u'(x)$ \\ 
\multicolumn{4}{r}{\text{innere Ableitung mal äußere Ableitung}}\\
\textbf{Bestimmtes Ingeral} &  &  $\Rightarrow$ &  $\int\limits_{a}^b f(x) \, \mathrm{d}x = F(b) - F(a)$
\end{tabularx}
\label{tab:ableitungsregeln}
\caption{Ableitungsregeln und Vorgehen}
\end{table}
 
\section{Wichtige Integrale}
\begin{align*}
\int f(x)\cdot g(x) \mathrm{d}x & =f(x)g(x)-\int f'(x)g(x)\mathrm{d}x
\end{align*}
Beispiel:
\begin{align*}
\int t\cdot e^{-j2\pi ft}\mathrm{d}t & =\frac{1}{-j2\pi f}\cdot\int t\cdot e^{-j2\pi ft}\mathrm{d}t=\frac{1}{-j2\pi f}\left[\left[t\cdot e^{-j2\pi ft}\right]-\int 1\cdot e^{-j2\pi ft}\mathrm{d}t\right]
\end{align*}

\begin{table}[!ht]
\centering
\begin{tabular}{l | c | c}
	$f(x)$ & $f'(x)$ & $\int f(x)\, \mathrm{d}x$ 	\\ 
	\toprule
	$x^n$ & $n \cdot x^{n-1}$ & $\begin{cases} \frac{1}{n+1} \cdot x^{n+1} + C  & \text{, für } n\neq -1 \\ ln|x| + C  &  \text{, für }n=-1 \end{cases}$ \\
	\midrule
	$e^x$ & $e^x$ & $e^x + C$ \\
	\midrule
	$e^{ax}$ &  $ae^{ax}$ & $\frac1a \cdot e^{ax}$ \\
	\midrule
	$\ln x$	& $\frac1x$ & $x(\ln((x)-1)) + C$ \\
	\midrule
	$\sin x$ & $\cos x$ & $-\cos(x) +C $\\
	\midrule
	$\cos x$ & $-\sin x$ & $\sin x$ \\
	\midrule
	$\frac{A}{(x-x_0)^n}$ & $\frac{-An(x-x_0)^{n-1}}{(x-x_0)^{2n}}$ &  $\begin{cases} A \cdot \ln |x-x_0| + C  &  \text{, für }n=1 \\ A \cdot \frac{1}{-n+1} \cdot \frac{1}{(x-x_0)^{n-1}}  & \text{, für}n>1	\end{cases} $\\
	\midrule
	$\e^{at}\sin(bt)$ & $a\e^{at}\sin(bt)+b\e^{at}\cos(bt)$ & $\frac{\e^{at}}{a^2+b^2}(a\sin(bt)-b\cos(bt))$\\
	\midrule
	$t\e^{at}$ & $\e^{at}+at\e^{at}$ & $\e^{at}(t-1)$
\end{tabular}
\caption{Wichtige Integrale und Ableitungen}
\end{table}

\section{Partialbruchzerlegung}
Bedingung: Nennergrad > Zählergrad, sonst ist eine Polynomdivision nötig.

Ansatz am Beispiel:
\begin{equation}
	\frac{3x^2-1}{(x-2)^2(x+3)x^3} = \frac{A}{x-2} + \frac{B}{(x-2)^2} + \frac{C}{x+3} + \frac{D}{x} + \frac{E}{x^2} + \frac{F}{x^3}
\end{equation}

Wenn der Nenner komplexe Nullstellen besitzt ist eine andere Zerlegung nötig, auch hier wieder ein Ansatz am Beispiel:
\begin{equation}
	f(x)=\frac{e^{-2x}}{(x+1)\cdot(x^2-2x+2)}=\frac{A}{x+1}+\frac{Bx+C}{x^2-2x+2}
\end{equation}

Jetzt werden A, B,$\dots$, F mit Hilfe von Nullstellen und Koeffizientenvergleich ermittelt. 
Sobald die Koeffizienten ermittelt sind kann die Rücktransformation erfolgen, die Regeln hierzu siehe Kapitel Laplace-Transformation bzw. z-Transformation.

\section{Quadratische Ergänzung}

\begin{align}
	a\cdot x^2 + b\cdot x +c\\
	\text{umformen zu}\\
	a \cdot \left( x + \left( \frac{b}{2\cdot a} \right)\right)^2 + c -a \cdot\left(\frac{b}{2\cdot a}\right)^2
\end{align}

\chapter{Fourier-Transformation}
\section{Grundlagen}
\subsection{Heaviside-Funktion $\sigma(t)$}
\subsubsection{Definition}
\begin{align}
	\sigma(t) = \begin{cases}
				1, & \text{für } t > 0\\
				0, & \text{für } t < 0\\
			\end{cases}\\
	\sigma(0) = \text{beliebig} \nonumber
\end{align}

\subsubsection{Verwendung}
\begin{align}
	&\text{Einschalten}\\\nonumber
	&f(t)\cdot\sigma(t-t_0)\\
	&\text{Ausschalten}\\\nonumber
	&f(t)\cdot\left[1-\sigma(t-t_0)\right]\\
	&\text{Ein- und Ausschalten}\\\nonumber
	&f(t)\cdot\left[\sigma(t-t_0) - \sigma(t-t_1)\right]
\end{align}

\subsection{$\delta$-Distribution}

\subsubsection{Definition}
\begin{align}
	\delta(t) = \begin{cases}
				0, & \text{für } t \neq 0\\
				\infty, & \text{für } t = 0\\
			\end{cases}\\
	\int_{-\infty}^\infty\!\delta(t)\,\mathrm{d}t = 1 \nonumber
\end{align}

\subsubsection{Zusammenhang $\sigma(t)$ und $\delta(t)$}
\begin{equation}
	\dot{\sigma}(t) = \delta(t)
\end{equation}

\section{Grundlagen Fourierreihen}
\subsection[Fourierreihe reell mit Periode $P=2\pi$]{Reelle Darstellung der Fourierreihe mit Periode $P=2 \pi$}
\begin{align*}
\begin{array}{rcll}
f(x) & = & \frac{a_0}{2} + \sum_{k=1}^{\infty} \; (a_k \cos(k x) + b_k \sin(k x)) & {}\\
\mbox{mit} & {} & {} & {}\\
	a_k & = & \frac{1}{\pi} \int_{-\pi}^{\pi} \; f(x) \cos(k x) \mathrm{d}x & \qquad k \in \{1,2,3, \ldots \}\\ 
	b_k & = & \frac{1}{\pi} \int_{-\pi}^{\pi} \; f(x) \sin(k x) \mathrm{d}x & \qquad k \in \{1,2,3, \ldots \}\\ 
	\frac{a_0}{2} & = & \frac{1}{2 \pi} \int_{-\pi}^{\pi} \; f(x) \mathrm{d}x
\end{array}
\end{align*}
\subsection[Fourierreihe reell mit Periode $T=\nicefrac{2\pi}{\omega}$]{Reelle Darstellung der Fourierreihe mit allgemeiner Periode $T=\nicefrac{2 \pi}{\omega}$}
\begin{align*}
\begin{array}{rcll}
f(t) & = & \frac{a_0}{2} + \sum_{k=1}^{\infty} \; (a_k \cos(k \omega t) + b_k \sin(k \omega t))  & {}\\ 
\mbox{mit} & {} & {} & {}\\
	a_k & = & \frac{2}{T} \int_{-T/2}^{T/2} \; f(t) \cos(k \omega t) \mathrm{d}t & \qquad k \in \{1,2,3, \ldots \}\\ 
	b_k & = & \frac{2}{T} \int_{-T/2}^{T/2} \; f(t) \sin(k \omega t) \mathrm{d}t & \qquad k \in \{1,2,3, \ldots \}\\ 
	\frac{a_0}{2} & = & \frac{1}{T} \int_{-T/2}^{T/2} \; f(t) \mathrm{d}t
\end{array}
\end{align*}
\subsection{Eigenschaften der reellen Fourierreihe}
\begin{itemize}
\item
$f(t)$ achsensymmetrisch zur $y$-Achse (gerade) $\Rightarrow$ $b_k = 0$, für alle $k \in \{1,2,3, \ldots \}$\\ 
$f(t)$ punktsymmetrisch zum Ursprung (ungerade) $\Rightarrow$ $a_k = 0$, für alle $k \in \{0,1,2,3, \ldots \}$\\
$f(t)$ punktsymmetrisch zu einem Punkt auf der $y$-Achse $\Rightarrow$ $a_k = 0$ für alle $k \in \{1,2,3, \ldots \}$, aber $a_0 / 2 \ne 0$. 
\item
	Klingen die \glqq{}langsamsten\grqq{} Fourierkoeffizienten ab wie $\nicefrac{1}{k^2}$, so ist $f(t)$ stetig.\\
	Klingen die \glqq{}langsamsten\grqq{} Fourierkoeffizienten ab wie $\nicefrac{1}{k}$, so hat $f(t)$ Sprungstellen.\\
(Diese Eigenschaften gelten natürlich auch für die komplexe Darstellung der Fourierreihen.)
\end{itemize}

\subsection[Fourierreihe komplex mit Periode $P=2\pi$]{Komplexe Darstellung der Fourierreihe mit Periode $P=2 \pi$}
\begin{align*}
\begin{array}{rcll}
f(x) & = & \sum_{k=-\infty}^{\infty} \; (c_k \cdot \e^{jkx}) & {}\\ 
\mbox{mit} & {} & {} & {}\\
	c_k & = & \frac{1}{2 \pi} \int_{-\pi}^{\pi} \; f(x) \cdot \e^{-jkx} \mathrm{d}x & \qquad k = 0,\pm 1,\pm 2,\pm 3, \ldots\\ 
	c_0 & = & \frac{1}{2 \pi} \int_{-\pi}^{\pi} \; f(x) \mathrm{d}x
\end{array}
\end{align*}
\subsection[Fourierreihe komplex mit Periode $T=\nicefrac{2\pi}{\omega}$]{Komplexe Darstellung der Fourierreihe mit allgemeiner Periode $T=\nicefrac{2 \pi}{\omega}$}
\begin{align*}
\begin{array}{rcll}
f(t) & = & \sum_{k=-\infty}^{\infty} \; (c_k \cdot \e^{jk \omega t}) & {}\\ 
\mbox{mit} & {} & {} & {}\\
	c_k & = & \frac{1}{T} \int_{-T/2}^{T/2} \; f(t) \cdot \e^{-jk \omega t} \mathrm{d}t & \qquad k = 0,\pm 1,\pm 2,\pm 3, \ldots\\ 
	c_0 & = & \frac{1}{T} \int_{-T/2}^{T/2} \; f(t) \mathrm{d}t
\end{array}
\end{align*}

\newpage

\noindent
\subsection[Umrechnung reelle und komplexe Koeffizienten]{Umrechnungsformeln zwischen den reellen und den komplexen Koeffizienten}
\begin{align*}
\begin{array}{rcl}
c_k & = & \frac{1}{2} ( a_k - j b_k)\\
c_{-k} = c^{*}_{k} & = & \frac{1}{2} ( a_k + j b_k)\\
c_0 & = & \frac{a_0}{2}\\
a_k & = & 2 \Re(c_k)\\
b_k & = & -2 \Im(c_k)
\end{array}
\end{align*}
\subsection[Amplituden-, Phasenspektrum, Fourierkoeffizienten]{Beziehungen zwischen Amplituden-, Phasenspektrum und den Fourierkoeffizienten}

\begin{align*}
\begin{array}{rcl}
A_k & = & \sqrt{a_{k}^{2} + b_{k}^{2}}\\
\tan(\varphi_k) & = & - \frac{b_k}{a_k}\\
|c_k| = |c_{-k}| & = & \frac{1}{2} A_k\\
\varphi_k & = & \arg( c_k)
\end{array}
\end{align*}


\section{Regeln zur Fouriertransformation}
\subsection{Fouriertransformation}
\begin{align*}
\begin{array}{lrcl}
	S(f) & = & \int_{-\infty}^{\infty} \, s(t) \cdot \e^{-j 2 \pi ft} \mathrm{d}t\\ 
	s(t) & = & \int_{-\infty}^{\infty} \, S(f) \cdot \e^{j 2 \pi ft} \mathrm{d}f\\ 
\end{array}
\end{align*}

\subsection{Rechenregeln Fouriertransformation}

\begin{align*}
\begin{array}{lrcl}
	\mbox{} & s(t) & \laplace & S(f)\\
\mbox{Linearität:} & a_1 s_1(t) + a_2 s_2(t) & \laplace & a_1 S_1(f) + a_2 S_2(f) \\
\mbox{Zeitverschiebungssatz:} & s(t-t_0) & \laplace & S(f) \cdot \e^{-j 2 \pi f t_0}\\
\mbox{Vertauschungssatz:} & S(t) & \laplace & s(-f) \\
\mbox{Symmetrie:} & s(-t) & \laplace & S(-f) \\
\mbox{Frequenzverschiebungssatz:} & s(t) \cdot \e^{j 2 \pi f_0 t}  & \laplace & S(f-f_0) \\
\mbox{Ähnlichkeitssatz:} & s(at) & \laplace & \frac{1}{|a|} \cdot S(\frac{f}{a})\\
\mbox{Differentiation im Zeitbereich:} & \frac{d^n s}{\mathrm{d}t^n} & \laplace & (j 2 \pi f)^n \cdot S(f)\\
\mbox{Differentiation im Frequenzbereich:} & (-j 2 \pi t)^n \cdot s(t) & \laplace & \frac{d^n S}{\mathrm{d}f^n}\\
\mbox{Multiplikationssatz:} & t^n \cdot s(t) & \laplace & \frac{j^n}{(2 \pi)^n} \cdot \frac{d^n S}{\mathrm{d}f^n}\\
\mbox{1. Faltungssatz:} & s_1(t) \star s_2(t) & \laplace & S_1(f) \cdot S_2(f)\\
\mbox{2. Faltungssatz:} & s_1(t) \cdot s_2(t) & \laplace & S_1(f) \star S_2(f)\\
\mbox{Integrationssatz:} & \int_{-\infty}^t s(\tau) d \tau & \laplace & \frac{1}{j 2 \pi f} \cdot S(f) + \frac{1}{2} S(0) \cdot \delta(f)\\
\mbox{Parsevalsche Gleichung:} & \int_{-\infty}^{\infty} s^2(t) \mathrm{d}t & = & \int_{-\infty}^{\infty} |S(f)|^2 \mathrm{d}f\\
\end{array}
\end{align*}

\subsection{Korrespondenzen Fouriertransformation}
\begin{align*}
\begin{array}{rcl}
\delta(t) & \laplace & 1\\
1 & \laplace & \delta(f)\\
\sigma(t) & \laplace & \frac{1}{j 2 \pi f} + \frac{1}{2} \delta(f)\\
\end{array}
\end{align*}

\subsection{Eigenschaften der Spektraldichte $S(f)$}
$\mbox{}$\\
Die Spektraldichte $S(f)$ einer geraden Zeitfunktion $s(t)$ ist reell und gerade.\\
Die Spektraldichte $S(f)$ einer ungeraden Zeitfunktion $s(t)$ ist rein imaginär; der Imaginärteil $I(f)$ dieser Spektraldichte ist ungerade.
$\mbox{}$\\

\subsection[Fourierkoeffizienten, Fourierintegral, periodische Funktion]{Beziehung zwischen Fourierkoeffizienten und Fourierintegral beim Übergang zu einer periodischen Funktion}
\begin{align*}
\begin{array}{rcl}
c_k & = & \frac{1}{T} \cdot S_0(k \cdot f_0)\\
	c_0 & = & \frac{1}{T} \int_{-T/2}^{T/2} \, f(t) \mathrm{d}t
\end{array}
\end{align*}


\subsection{Faltung}
\subsubsection{Definition}
\begin{equation}
	f(t) * g(t) = \int_{-\infty}^{\infty}\!f(\tau) \cdot g(t-\tau) \, \mathrm{d}\tau 
\end{equation}

\clearpage

\subsubsection{Berechnung}
\begin{table}[h]
	\begin{tabular}{ll}
		Allgemeines Vorgehen & Beispiel $f(t) * g(t) $\\
				     &$f(t)=\sigma(t)\cdot\e^{-t}$\\
		       & $g(t) = \sigma(t+1) - \sigma(t-1)$\\
		\toprule
		1. Reihenfolge festlegen (mehr $\sigma \Rightarrow$ zuerst) & $g(t)*f(t)$\\
									    &\\
		2. Definition der Faltung übernehmen & $\int_{-\infty}^{\infty}\!g(\tau)\cdot f(t-\tau)\,\mathrm{d}\tau$\\
						      &\\
		3. Fallunterscheidung nach $\tau$ (wann werden die $\sigma=1$?) & $0 < \tau < 1$ \&  $\tau < t$\\
										&\\
		4. Fallunterscheidung nach t & $\begin{cases} 	1. & t < 0 \Rightarrow \int\dots = 0\\ 
								2. & t \in \left[0;1\right] \Rightarrow \int_0^t\e^{-(t-\tau)}\mathrm{d}\tau\\
								3. & t > 1 \Rightarrow \int_0^1\e^{-(t-\tau)}\mathrm{d}\tau
						\end{cases}$\\
						&\\
		5. Integration & Integralbildung\\
							   &\\
		6. Zusammenfassung & $f(t) * g(t) = \begin{cases}
							0, & t < 0\\
						1-\e^{-t}, & 0 < t < 1\\
						\e^{-t}(e-1), & t > 1
						\end{cases}$\\
	\end{tabular}
	\caption{Vorgehen Faltung}
	\label{tab:vorgehen_faltung}
\end{table}

\clearpage


\chapter{Laplace-Transformation}
\section{Die Laplacetransformation}
\begin{align*}
\begin{array}{rcl}
F(p) & = & \int_{0}^{\infty} \, f(t) \cdot \e^{-pt} dt\\ 
f(t) & = & \frac{1}{j 2 \pi} \int_{p=\alpha-j \cdot \infty}^{p=\alpha +j \cdot \infty} \, F(p) \cdot \e^{pt} dp\\ 
\end{array}
\end{align*}

\section{Regeln Laplacetransformation}

\begin{align*}
\begin{array}{lrcl}
	\mbox{} & f(t) & \laplace & F(p)\\
\mbox{Linearität:} & \sigma(t) \cdot (a_1 f_1(t) + a_2 f_2(t)) & \laplace & a_1 F_1(p) + a_2 F_2(p) \\
\mbox{Ähnlichkeitssatz:} & \sigma(t) \cdot f(at); \quad a>0  & \laplace & \frac{1}{a} \cdot F(\frac{p}{a}) \\
\mbox{Verschiebungssatz:} & f(t-a) \cdot \sigma(t-a); \quad a>0  & \laplace & \e^{-ap} \cdot F(p)\\
\mbox{Dämpfungssatz:} & \sigma(t) \cdot \e^{-at} \cdot f(t) & \laplace & F(p+a)\\
\mbox{1. Differentiationssatz:} & \sigma(t) \cdot f'(t) & \laplace & p \cdot F(p) - f(0+)\\
\mbox{2. Differentiationssatz:} & \sigma(t) \cdot -t \cdot f(t) & \laplace & \frac{d}{dp} F(p)\\
\mbox{Integrationssatz:} & \sigma(t) \cdot \int_0^t f(\tau) d \tau & \laplace & \frac{1}{p} \cdot F(p)\\
\mbox{Faltungssatz:} & \sigma(t) \cdot f_1(t) \star f_2(t) & \laplace & F_1(p) \cdot F_2(p)
%\mbox{Endwertsätze:} & \lim_{t \to \infty} f(t) & = & \lim_{p \to 0} F(p)\\
%\mbox{} & \lim_{t \to 0} f(t) & = & \lim_{p \to \infty} F(p)\\
\end{array}
\end{align*}

\section{Höhere Ableitungen}
\begin{align*}
\begin{array}{lcl}
f'(t) & \laplace & p \cdot F(p) - f(0+)\\
f''(t) & \laplace & p^2 \cdot F(p) -p \cdot f(0+) - f'(0+)\\
f'''(t) & \laplace & p^3 \cdot F(p) - p^2 \cdot f(0+) - p \cdot f'(0+) - f''(0+)\\
\mbox{usw.} & {} & {}
\end{array}
\end{align*}

\section{Korrespondenzen Laplacetransformation}
\begin{align*}
\begin{array}{rcll}
	\sigma(t) \cdot 1 & \laplace & \frac{1}{p}; & \quad \mbox{Re}(p) > 0\\
	\sigma(t) \cdot t^n & \laplace & \frac{n!}{p^{n+1}}; & \quad \mbox{Re}(p) > 0\\
	\sigma(t) \cdot \frac{1}{p^n} & \Laplace & \frac{t^{n-1}}{(n-1)!}; & \quad \mbox{Re}(p) > 0\\
	\sigma(t) \cdot \e^{at} & \laplace & \frac{1}{p-a}; & \quad \mbox{Re}(p) > \mbox{Re}(a)\\
	\sigma(t) \cdot \delta(t) & \laplace & 1 & {}\\
	\sigma(t) \cdot \cos(\omega t) & \laplace & \frac{p}{p^2 + \omega^2} & {}\\
	\sigma(t) \cdot \sin(\omega t) & \laplace & \frac{\omega}{p^2 + \omega^2} & {}
\end{array}
\end{align*}

\section{Weitere nützliche Formeln und Hinweise}
\begin{align*}
\begin{array}{lcl}
\e^{j \varphi} + \e^{-j \varphi} & = & 2 \cos \varphi\\ 
\e^{j \varphi} - \e^{-j \varphi} & = & 2 j \sin \varphi\\ 
\end{array}
\end{align*}

Hinweise:\\
\begin{itemize}
	\item Finden sich im Nenner einer Funktion komplexe Nullstellen enthält deren Laplace-Transformierte auch trigonometrische Funktionen
\end{itemize}


\section{Differentialgleichungen mit Faltung lösen}
\begin{itemize}
	\item Linke Seite transformieren, rechte Seite \glqq{}$\mathcal{L}(r(t))$\grqq{}
	\item Auflösen nach $X(p)$
	\item Bestimmen der Laplace-Inversen mit p-Anteilen, $\sigma$ nicht vergessen
	\item Faltung anwenden
\end{itemize}


\chapter{z-Transformation}
\section{Regeln $z$-Transformation} 

\begin{align*}
\{ a_k \} \;\; \laplace \;\; \sum_{k=0}^{\infty} a_k \cdot z^{-k}
\end{align*}

\begin{align*}
	\mbox{Verschiebungssatz:}\qquad \{ f_{k-n} \} \;\; \laplace \;\; z^{-n} \cdot F(z)
\end{align*}

\subsection{Korrespondenzen $z$-Transformation}
\begin{align*}
\begin{array}{rcc}
	\{ \sigma_k \} & \laplace & \frac{z}{z-1}\\
	\{ \delta_k \} & \laplace & 1\\
	\{ a^k \cdot \sigma_k \} & \laplace & \frac{z}{z-a}
\end{array}
\end{align*}

\subsection{Aus dem Verschiebungssatz folgende Korrespondenzen}
\begin{align*}
\begin{array}{rcll}
\{ \sigma_{k-1} \} & \laplace & \frac{1}{z-1}\\
\{ a^{k-1} \cdot \sigma_{k-1} \} & \laplace & \frac{1}{z-a}
\end{array}
\end{align*}


\section{Am Beispiel $y_n - a^2 y_{n-2} = \delta_n$}

\begin{align*}
	y_n - a^2\cdot y_{n-2} = \delta_n \laplace Y(z) - a^2 \frac{1}{z^2}Y(z) = 1\\
	\\
	Y(z)\left(1-\frac{a^2}{z^2}\right) = 1\\
	Y(z)\left(\frac{z^2-a^2}{z^2}\right) = 1\\
	Y(z) = \frac{1}{\frac{z^2-a^2}{z^2}} = \frac{z^2}{z^2-a^2}\\
	Y(z) = \frac{z^2-a^2+a^2}{z^2-a^2} = 1 + \frac{a^2}{z^2-a^2}\\
	\\
	\Rightarrow \text{Beginne Partialbruchzerlegung}\\
	\\
	\frac{a^2}{z^2-a^2} = \frac{A}{z-(-a)}+\frac{B}{z-a}\\
	\Rightarrow A=\frac{-a}{2}\\
	B=\frac{a}{2}\\
	\\
	\Leftarrow \text{Beende Partialbruchzerlegung}\\
	\\
	\Rightarrow Y(z) = 1-\frac{a}{2}\frac{1}{z-(-a)}+\frac{a}{2}\frac{1}{z-a}\\
	Y(z) \Laplace y_n = \delta - \frac{a}{2}(-a)^{n-1}\sigma_{n-1}+\frac{a}{2}a^{n-1}\sigma_{n-1}\\
	\\
	y_n = \delta_n+\frac{1}{2}(-a)^n\sigma_{n-1}+\frac{1}{2}a^n\sigma_{n-1}
\end{align*}

\chapter{Statistik}
\section{Häufigkeiten}
\begin{table}[h]
	\begin{tabular}{lll}
		Variable & Formel & Bedeutung\\
		\toprule
		$n_j$ & & absolute Häufigkeit\\
		$N$ & $\sum\limits_j n_j$ & Anzahl Messdaten\\
		$h_j$ & $\frac{n_j}{N}$ & relative Häufigkeit, $\sum\limits_j h_j = 1$\\
		$N_j$ & $\sum\limits_{k=1}^{j} n_k$ & absolute Summenhäufigkeit\\
		$H_j$ & $\sum\limits_{k=1}^{j} h_k$ & relative Summenhäufigkeit, $\sum\limits_j H_j = 1$\\
	\end{tabular}
\end{table}

\section{Klassische Wahrscheinlichkeit}
\begin{equation}
P(A) = \frac{|A|}{|\Omega|} = \frac{\text{Anzahl günstiger Fälle}}{\text{Anzahl möglicher Fälle}}
\end{equation}

\section{Permutationen}
Eine Anordnung von $n$ Elementen in einer bestimmten Reihenfolge heißt Permutation.

Für die Anzahl gilt:
\begin{equation}
	P(n) = n!
\end{equation}

Sind $n_1, n_2, n_3, \dots, n_k$ Elemente jeweils einander gleich:\\
\begin{equation}
	P(n; n_1,n_2,\dots,n_k) = \frac{n!}{n_1!\cdot n_2!\dots n_k!}
\end{equation}

\section{Kombinatorik}
\begin{table}[h]
	\begin{tabular}{l|ll}
				& Ohne Zurücklegen & Mit Zurücklegen \\
		\toprule
		geordnet	& $N=\frac{n!}{(n-k)!}$ & $N=n^k$\\
		\midrule
		ungeordnet	& $N=\binom{n}{k}=\frac{n!}{(n-k)!k!}$ & $N=\binom{n+k-1}{k}$
	\end{tabular}
	\caption{Entscheidungstabelle Wahrscheinlichkeiten}
\end{table}

\section{Variablen}
\begin{table}[h]
	\begin{tabular}{lll}
		Variable & Bedeutung & Beispiel\\
		\toprule
		$n$ & Größe der Grundgesamtheit & Kugeln in der Urne\\
		$k$ & Größe der Stichprobe & gezogene Kugeln\\
		$N$ & Anzahl der möglichen Kombinationen & Ergebnis
	\end{tabular}
	\caption{Variablen}
\end{table}

\section{Rechenregeln der Wahrscheinlichkeit}
\begin{align*}
	& P(\emptyset) = \frac{|\emptyset|}{|\Omega|}=0\\
	& P(\Omega) = \frac{|\Omega|}{|\Omega|}=1\\
	& \Rightarrow  0\leq P(A) \leq 1\\
	& P(A\cup B) = P(A)+P(B)-P(A\cap B)\\
	& P(A)+P(\bar{A}) =1\\
	& P(A\cup B) = P(A) + P(B) \Leftrightarrow A\cup B = \emptyset
\end{align*}

\section{Bedingte Wahrscheinlichkeit und Unabhängigkeit}
\begin{equation}
	P(A|B) = \frac{P(A\cap B)}{P(B)}
\end{equation}
\begin{equation}
	P(A\cap B)=P(A|B)\cdot P(B)
\end{equation}
\begin{equation}
	P(A\cap B)=P(B|A)\cdot P(A)
\end{equation}
$\Rightarrow$ Multiplikationssatz:
\begin{equation}
	P(A|B)\cdot P(B) = P(B|A)\cdot P(A)
\end{equation}
\begin{equation}
	P(A\cap B) = P(A)\cdot P(B) \Rightarrow \text{A unabhängig von B}
\end{equation}

\section{Totale Wahrscheinlichkeit}
\subsection{Voraussetzungen}
\begin{equation*}
	B_i \cap B_j = \emptyset\\
\end{equation*}
\begin{equation*}
	B_1 \cup \dots \cup B_n = \Omega
\end{equation*}
\begin{equation*}
	A \subseteq \Omega
\end{equation*}

\subsection{Formel}
\begin{equation}
	P(A) = P(A|B_1)\cdot P(B_1) + P(A|B_2) \cdot P(B_2) + \dots + P(A|B_n)\cdot P(B_n)
\end{equation}

\subsection{Pfadregeln}
Parallel verlaufende Pfade werden addiert, hintereinander folgende Pfade multipliziert.\\
Beispiel:\\
Maschine, 3 Bauteile, 2 davon parallel zueinander, danach das dritte. Die Ausfallwahrscheinlichkeiten der beiden ersten werden addiert, dann mit der Ausfallwahrscheinlichkeit des dritten multipliziert um die Ausfallwahrscheinlichkeit der gesamten Maschine zu berechnen.

\clearpage

\section{Bayes}
\subsection{Voraussetzungen}
\begin{equation*}
	B_i \cap B_j = \emptyset\\
\end{equation*}
\begin{equation*}
	B_1 \cup \dots \cup B_n = \Omega
\end{equation*}
\begin{equation*}
	A \subseteq \Omega
\end{equation*}

\subsection{Formel}
\begin{equation}
	P(B_k|A) = \frac{P(A|B_k)\cdot P(B_k)}{\sum\limits_{k=1}^n P(A|B_k) \cdot P(B_k)}
\end{equation}

\section{Zufallsvariable}
Sei $\omega \in \Omega$ ein Elementarereignis. Dann ist $X(\omega)$ Zufallsvariable. $X(\omega)$ hängt von der Fragestellung ab.

\section{Wahrscheinlichkeitsdichte}
\begin{equation}
	f(x) = P(X = x)
\end{equation}

\section{Diskrete ZV}
Die aufsummierte Wahrscheinlichkeitsdichte heißt Verteilungsfunktion $F(x)$\\
Erwartungswert $\mu$:
\begin{align}
	& \mu = E(X) = \sum_i x_i \cdot f(x_i)\\
	& E(nX) = nE(X)\\
	& E(X_1 + X_2) = E(X_1) + E(X_2)
\end{align}

Fairness:\\
Ist der Erwartungswert $E(X) = 0$ gilt beispielsweise ein Würfelspiel als fair.\\


Varianz $\sigma^2$:
\begin{align}
	& \sigma^2 = \mathrm{Var}(X)=\sum_i(x_i-\mu)^2\cdot f(x_i) = \sum_i x_i^2f(x_i)-\mu^2\\
	& \mathrm{Var}(nX) = n^2\mathrm{Var}(X)\\
	& \mathrm{Var}(nX_1+mX_2) = n^2\mathrm{Var}(X_1) + m^2\mathrm{Var}(X_2)
\end{align}

Standardabweichung $\sigma$:
\begin{equation}
	\sigma = \sqrt{\sigma^2}
\end{equation}

\clearpage

\section{Stetige ZV}
Erwartungswert $\mu$:
\begin{equation}
	\mu = \int_{-\infty}^{\infty}xf(x)\mathrm{d}x\\
\end{equation}
Varianz $\sigma^2$:
\begin{equation}
	\sigma^2 = \int_{-\infty}^{\infty}x^2f(x)\mathrm{d}x-\mu^2\\
\end{equation}
Standardabweichung $\sigma$:
\begin{equation}
	\sigma = \sqrt{\sigma^2}
\end{equation}
\begin{align}
	P(X\leq a) &= F(a)\\
	P(X > a) &= 1 - F(a)\\
	P(a < X < b) &= F(b) - F(a)
\end{align}

Außerdem gilt:
\begin{align}
	F(-\infty) &= 0\\
	F(\infty) &= 1\\
	\int_{-\infty}^\infty\!f(x)\, \mathrm{d}x &= 1
\end{align}

\clearpage

\section{Verteilungen}
\subsection{Binomialverteilung}
\large{Ziehen mit Zurücklegen, Bernoulli-Experiment}\\
Diskrete Dichte:
\begin{equation}
	P(X=k) = \binom{n}{k}p^k(1-p)^{n-k}
\end{equation}

Erwartungswert:
\begin{equation}
	\mu = E(X) = np
\end{equation}

Varianz:
\begin{equation}
	\text{Var}(X)=np(1-p)
\end{equation}

Approximation durch Poisson-Verteilung:
\begin{align*}
	& \text{Faustregel} \\
	& np \leq 10 \text{ und} \\
	& n \geq 1500p \Rightarrow Ps(\mu = np) 
\end{align*}

Approximation durch Normalverteilung:
\begin{align*}
	& \text{Faustregel} \\
	& np(1-p) > 9 \Rightarrow N(\mu = np; \sigma = \sqrt{np(1-p)} 
\end{align*}

Anwendungsbeispiel:\\
Würfelspiel, 3 Würfel, einer davon zeigt die gewünschte Augenzahl:\\
\begin{equation}
	\binom{3}{1}\cdot\left(\frac{1}{6}\right)^1\cdot\left(1-\frac{1}{6}\right)^2
\end{equation}

\begin{table}[h]
	\begin{tabular}{ll}
		Variable & Erklärung\\
		\toprule
		$n$ & Anzahl Würfel\\
		$k$ & Anzahl Treffer (günstiges Ereignis)\\
		$p$ & Wahrscheinlichkeit einer bestimmten Augenzahl
	\end{tabular}
\end{table}

\clearpage

\subsection{Hypergeometrische Verteilung}
\large{Ziehen ohne Zurücklegen}\\
Diskrete Dichte:
\begin{equation}
	P(X=k)=\frac{\binom{M}{k}\binom{N-M}{n-k}}{\binom{N}{n}}
\end{equation}

Erwartungswert:
\begin{equation}
	\mu = E(X)= n\frac{M}{N}
\end{equation}

Varianz:
\begin{equation}
	\mathrm{Var}(X)=n\frac{M}{N}\left(1-\frac{M}{N}\right)\left(\frac{N-n}{N-1}\right)
\end{equation}

Approximation durch Binomialverteilung
\begin{align*}
	& \text{Faustregel} \\
	& 0,1 < \frac{M}{N} < 0,9 \\
	& n < 0,05N, n >10 \Rightarrow B\left(n;p = \frac{M}{N}\right) 
\end{align*}

Approximation durch Poisson-Verteilung
\begin{align*}
	& \text{Faustregel} \\
	& \frac{M}{N}\leq 0,1 \text{ oder } \frac{M}{N}\geq 0,9 \\
	& n < 0,05N, n > 30 \Rightarrow Ps\left(\mu = n\frac{M}{N}\right) 
\end{align*}

Approximation durch Normalverteilung
\begin{align*}
	& \text{Faustregel}  \\
	& 0,1 < \frac{M}{N} < 0,9  \\
	& n < 0,05N, n > 30 \Rightarrow N\left(\mu = n\frac{M}{N}; \sigma = \sqrt{n\frac{M}{N}\left(1-\frac{M}{N}\right)\frac{N-n}{N-1}}\right) 
\end{align*}

Anwendungsbeispiel:\\
N Kugeln, M davon rot, n Kugeln werden gezogen, P ist die Wahrscheinlichkeit, dass von n Kugeln genau k rot sind.

\subsection{Poissonverteilung}
\large{Mehrstufiges Bernoulli-Experiment}\\
Diskrete Dichte:
\begin{equation}
	P(X=k)=\frac{\lambda^k}{k!}\e^{-\lambda}
\end{equation}

Erwartungswert:
\begin{equation}
	\mu = E(X) = \lambda
\end{equation}

Varianz:
\begin{equation}
	\mathrm{Var}(X)=\lambda
\end{equation}

Approximation durch Normalverteilung
\begin{align*}
	& \text{Faustregel} \\
	& \mu > 10 \Rightarrow N(\mu; \sigma = \sqrt{\mu}) 
\end{align*}

Anwendungsbeispiel:\\
Blitze pro Stunde pro $\mathrm{km}^2$, P ist die Wahrscheinlichkeit, dass ein Ereignis k-mal eintritt, wobei es im Mittel $\lambda$-mal eintritt.

\clearpage

\subsection{Normalverteilung}
Dichtefunktion der NV mit $\mu$ und $\sigma^2$:
\begin{equation}
	f(x) = \frac{1}{\sqrt{2\pi\sigma^2}}\e^{-\frac{1}{2}\frac{(x-\mu)^2}{\sigma^2}}
\end{equation}

Verteilungsfunktion:
\begin{equation}
	F(x) = \int_{-\infty}^xf(t)\mathrm{d}t
\end{equation}

Da $f(x)$ sich nicht integrieren lässt, lässt sich $F(x)$ nicht analytisch berechnen. Ausweichen auf die Standardnormalverteilung $\Phi(z)$ samt Tabelle.

Standardnormalverteilung hat $\mu=0$ und $\sigma^2=1$.\\
Dichtefunktion:
\begin{equation}
	\Phi(z) = \frac{1}{\sqrt{2\pi}}\e^{-\frac{1}{2}x^2}
\end{equation}

Berechnung von Wahrscheinlichkeiten mit der NV:\\
\begin{align}
	\begin{rcases}
		X \text{ sei } N(\mu, \sigma^2)\text{-verteilt} &\\
		Z \text{ sei } N(0,1)\text{-verteilt} &
	\end{rcases} \text{Umrechnungsformel: } z=\frac{x-\mu}{\sigma}
\end{align}

\begin{align}
	P(X < a) &= F(a) = \Phi\left(\frac{a-\mu}{\sigma}\right)\\
	P(X > a) &= 1-F(a) = 1-\Phi\left(\frac{a-\mu}{\sigma}\right)\\
	P(a < X < b) &= F(b) - F(a) = \Phi\left(\frac{b-\mu}{\sigma}\right) - \Phi\left(\frac{a-\mu}{\sigma}\right)
\end{align}

\subsection{Quantile}
Quantile der Normalverteilung:\\
$x_{0,95}$ heißt 95\%-Quantil.\\
Es gilt: $P(X < x_{0,95}) = F(X_{0,95}) = 0,95$.\\

Quantile der Standardnormalverteilung:\\
$z_{0,8}, z_{0,95}\dots$ siehe Tabelle.\\

Umrechnung:\\
$x_{0,95} = \mu +\sigma\cdot z_{0,95}$\\

Zu beachten: $\alpha=$Signifikanzniveau (Irrtumswahrscheinlichkeit).\\
Schreibweise: $x_{1-\alpha}=\mu+\sigma\cdot z_{1-\alpha}$

\section{Zufallsstreubereiche}
95\%-ZSB: Intervall, in dem $X$ mit 95\%-iger Wahrscheinlichkeit liegt.

\begin{itemize}
	\item zweiseitiger ZSB: $\left[\mu-z_{1-\frac{\alpha}{2}}\sigma; \mu+z_{1-\frac{\alpha}{2}}\sigma\right]$
	\item einseitiger, nach oben beschränkter ZSB: $(-\infty;\mu+z_{1-\alpha}\sigma]$
	\item einseitiger, nach unten beschränkter ZSB: $[\mu-z_{1-\alpha}\sigma; \infty)$
\end{itemize}

\section{Zufallsstreubereich für $\bar{X}$}

\begin{itemize}
	\item zweiseitiger ZSB: $\left[\mu-z_{1-\frac{\alpha}{2}}\frac{\sigma}{\sqrt{n}}; \mu+z_{1-\frac{\alpha}{2}}\frac{\sigma}{\sqrt{n}}\right]$
	\item einseitiger, nach oben beschränkter ZSB: $(-\infty;\mu+z_{1-\alpha}\frac{\sigma}{\sqrt{n}}]$
	\item einseitiger, nach unten beschränkter ZSB: $[\mu-z_{1-\alpha}\frac{\sigma}{\sqrt{n}}; \infty)$
\end{itemize}

\section{Schließende Statistik}
\subsection{Punktschätzer}
Es liegt eine Messreihe $X_1, X_2, \dots, X_n$ vor.

\begin{table}[ht!]
	\begin{tabular}{ll}
		Unbekannter Parameter & Punktschätzer\\
		\toprule
		Wahrscheinlichkeit $p$ & $\hat{p}=\frac{k}{n}\dots$ relative Häufigkeit\\
		Erwartungswert $\mu$ & $\hat{\mu}=\bar{x}\dots$ arithmetisches Mittel\\
		Varianz $\sigma^2$ & $\hat{\sigma}^2 = s^2\dots$ empirische Varianz\\
		Standardabweichung $\sigma$ & $\hat{\sigma} = s\dots$ empirische Standardabweichung
	\end{tabular}
	\caption{Punktschätzer}
\end{table}

\subsection{Statistische Tests}
Ablauf:
\begin{itemize}
	\item Formulierung einer Nullhypothese $H_0$ und der zugehörigen Alternativhypothese $H_1$
	\item Wahl des Signifikanzniveaus $\alpha$
	\item ZSB wird nach Tabelle je nach Testart berechnet
	\item Testentscheidung, wenn $\bar{x}\notin\mathrm{ZSB}\Rightarrow H_0$ wird verworfen $\Rightarrow H_1$ gilt als bestätigt, mit Irrtumswahrscheinlichkeit $\alpha$, wenn $\bar{x}\in\mathrm{ZSB}\Rightarrow$ keine Aussage möglich
\end{itemize}

\subsection{Konfidenzintervalle}
Konfidenzintervalle bzw. Vertrauensbereiche sind Intervalle, in denen $\mu$ (Fälle A \& B), $\mu_1-\mu_2$ (Fall C) bzw. $p$ (Fall D) mit der Wahrscheinlichkeit $1-\alpha$ liegt. Es gibt zweiseitige, einseitig unten und einseitig oben beschränkte Konfidenzintervalle.

\clearpage

\section{Beschreibende Statistik}

\subsection{Arithmetisches Mittel}
\begin{align*}
\overline{x} = \frac{1}{n} \cdot \sum_{i=1}^{n} x_i
\end{align*}

\subsection{Empirische Varianz}
\begin{align*}
	s^2 &= \frac{1}{n-1} \cdot \sum_{i=1}^{n} \left(x_i - \overline{x}\right)^2\\
	    &= \frac{1}{n-1} \cdot \left(\left(\sum_{i=1}^{n} x_i^2 \right) - n \cdot \overline{x}^2\right)
\end{align*}

\subsection{Empirische Standardabweichung}
\begin{align*}
	s = \sqrt{s^2} = \sqrt{\frac{1}{n-1} \cdot \sum_{i=1}^{n} (x_i - \overline{x})^2} = \sqrt{\frac{1}{n-1} \cdot ((\sum_{i=1}^{n} x_i^2) - n \cdot \overline{x}^2)}
\end{align*}

\subsection{Empirische Kovarianz}
\begin{align*}
	s_{xy} &= \frac{1}{n-1} \cdot \sum_{i=1}^{n} (x_i-\overline{x})(y_i-\overline{y})\\
	&= \frac{1}{n-1} \cdot \left(\sum_{i=1}^{n} (x_i \cdot y_i) - n \cdot \overline{x} \cdot \overline{y}\right)\\
\end{align*}

Anwendungsbeispiel: Kovarianz von Gewicht und Haltbarkeit von Werkzeugen\\
Ist das Vorzeichen negativ nimmt die Haltbarkeit mit steigendem Gewicht des Werkzeugs ab, ist es positiv nimmt sie zu.

\subsection{Empirischer Korrelationskoeffizient}
\begin{align*}
	r_{xy} = \frac{s_{xy}}{s_x \cdot s_y}\\
	-1 \leq r_{xy} \leq 1
\end{align*}

Für $|r_{xy}|\rightarrow 1$ gilt starke Abhängigkeit von $x$ und $y$, für $|r_{xy}|\rightarrow 0$ ist die Abhängigkeit schwach.

\subsection{Regressionsgerade}
\begin{align*}
	Q(m,b) &= \sum_{i=1}^{n}\left[y_i-(mx_i+b)\right]^2=\text{Min!}\\
	m &= \frac{s_{xy}}{s_x^2}\\
	b &= \bar{y}-m\cdot\bar{x}
\end{align*}

Alternative Formel:

\begin{align*}
	y &= ax+b\\
\\
	a &= \frac{n \sum\limits_{i=1}^{n} x_i y_i - \left( \sum\limits_{i=1}^{n} x_i \right) \left( \sum\limits_{i=1}^{n} y_i \right)}
		{n \sum\limits_{i=1}^{n} x_i^2 - \left( \sum\limits_{i=1}^{n} x_i \right)^2}\\
\\
	b &= \frac{\left( \sum\limits_{i=1}^{n} x_i^2 \right) \left( \sum\limits_{i=1}^{n} y_i \right)-
		\left( \sum\limits_{i=1}^{n} x_i \right) \left( \sum\limits_{i=1}^{n} x_i y_i \right)}
		{n \sum\limits_{i=1}^{n} x_i^2- \left( \sum\limits_{i=1}^{n} x_i \right)^2}
\end{align*}

\subsection{Median}
Der Median ist die exakte Mitte einer geordneten Messreihe.

\begin{equation*}
	\tilde{x} = x_{0,5}
\end{equation*}


\clearpage

\section{Statistische Tests}

\subsection{Test für den Parameter $\mu$ einer Normalverteilung bei bekannter Varianz (Gauß-Test)}
Stichprobe $x_1, \ldots , x_n$\\
Messwerte sind Realisierungen von $n$ unabhängigen $N(\mu, \sigma^2)$ verteilten Zufallsvariablen mit unbekanntem Erwartungswert $\mu$, aber bekannter Varianz $\sigma^2$.\\
Signifikanzniveau $\alpha$\\

\begin{tabular}{llll}
	$H_0$ 		& $H_1$ 		& Zufallsstreubereich (ZSB) für $\overline{x}$, falls $H_0$ zutrifft & $H_0$ verwerfen falls\\
	\toprule
	$\mu = \mu_0$ 	& $\mu \ne \mu_0$ 	& $[\mu_0 - z_{1-\alpha /2} \cdot \frac{\sigma}{\sqrt{n}} ; \;
	\mu_0 + z_{1-\alpha /2} \cdot \frac{\sigma}{\sqrt{n}}]$  & $\overline{x} \not\in$ ZSB \\
	$\mu \ge \mu_0$ & $\mu < \mu_0$ & $[\mu_0 - z_{1-\alpha} \cdot \frac{\sigma}{\sqrt{n}} ; \;
	\infty)$ & $\overline{x} \not\in$ ZSB \\
	$\mu \le \mu_0$ & $\mu > \mu_0$ & $(-\infty ; \;
	\mu_0 + z_{1-\alpha} \cdot \frac{\sigma}{\sqrt{n}}]$  & $\overline{x} \not\in$ ZSB 
\end{tabular}


\subsection{Test für den Parameter $\mu$ einer Normalverteilung bei unbekannter Varianz ($t$-Test)}

Stichprobe $x_1, \ldots , x_n$\\
Messwerte sind Realisierungen von $n$ unabhängigen $N(\mu, \sigma^2)$ verteilten Zufallsvariablen mit unbekanntem Erwartungswert $\mu$ und unbekannter Varianz $\sigma^2$.\\
Die unbekannte Standardabweichung $\sigma$ wird durch die empirische Standardabweichung $s$ aus der Stichprobe geschätzt.\\
Signifikanzniveau $\alpha$\\


\begin{tabular}{llll}
	$H_0$ 		& $H_1$ 	& Zufallsstreubereich (ZSB) für $\overline{x}$, falls $H_0$ zutrifft & $H_0$ verwerfen falls\\
	\toprule
	$\mu = \mu_0$ 	& $\mu \ne \mu_0$ & $[\mu_0 - t_{n-1;1-\alpha /2} \cdot \frac{s}{\sqrt{n}} ; \;
	\mu_0 + t_{n-1;1-\alpha /2} \cdot \frac{s}{\sqrt{n}}]$  & $\overline{x} \not\in$ ZSB \\ 
	$\mu \ge \mu_0$ & $\mu < \mu_0$ & $[\mu_0 - t_{n-1;1-\alpha} \cdot \frac{s}{\sqrt{n}} ; \;
	\infty)$ & $\overline{x} \not\in$ ZSB \\
	$\mu \le \mu_0$ & $\mu > \mu_0$ & $(-\infty ; \;
	\mu_0 + t_{n-1;1-\alpha} \cdot \frac{s}{\sqrt{n}}]$  & $\overline{x} \not\in$ ZSB 
\end{tabular}

\clearpage

\subsection{Vergleich der Parameter $\mu_1$ und $\mu_2$ zweier Normalverteilungen (Zweistichproben-$t$-Test)}

Zwei Stichproben $x_1, \ldots , x_m$ und $y_1, \ldots , y_n$\\
Die $m$ Messwerte $x_1, \ldots , x_m$ sind Realisierungen von $N(\mu_1, \sigma^2)$ verteilten Zufallsvariablen; die $n$ Messwerte $y_1, \ldots , y_n$ sind Realisierungen von $N(\mu_2, \sigma^2)$ verteilten Zufallsvariablen. Alle Zufallsvariablen sind voneinander unabhängig mit der gleichen unbekannten Varianz $\sigma^2$.\\
Aus den beiden empirischen Varianzen $s_x^2$ der ersten Stichprobe und $s_y^2$ der zweiten Stichprobe muss zunächst die folgende Hilfsgröße berechnet werden:
\begin{align*}
s_d = \sqrt{(m-1) s_{x}^2 + (n-1) s_{y}^2} \cdot \sqrt{\frac{m+n}{m \cdot n \cdot (m+n-2)}}
\end{align*}
Signifikanzniveau $\alpha$\\

\begin{tabular}{llll}
	$H_0$ 		& $H_1$ 	& Zufallsstreubereich (ZSB) für $\overline{x} - \overline{y}$, falls $H_0$ zutrifft & $H_0$ verwerfen falls \\
	\toprule
	$\mu_1 = \mu_2$ & $\mu_1 \ne \mu_2$ & $[- t_{m+n-2;1-\alpha /2} \cdot s_d ; \;
	t_{m+n-2;1-\alpha /2} \cdot s_d]$  & $\overline{x}-\overline{y} \not\in$ ZSB \\
	$\mu_1 \ge \mu_2$ & $\mu_1 < \mu_2$ & $[- t_{m+n-2;1-\alpha} \cdot s_d ; \; \infty)$  & $\overline{x}-\overline{y} \not\in$ ZSB \\
	$\mu_1 \le \mu_2$ & $\mu_1 > \mu_2$ & $(-\infty ; \;
	t_{m+n-2;1-\alpha} \cdot s_d]$  & $\overline{x}-\overline{y} \not\in$ ZSB
\end{tabular}

\subsection{Test für den Parameter $p$ einer Binomialverteilung}

Hier wird ein Test für ein Ereignis, das mit unbekannter Wahrscheinlichkeit $p$ auftritt, durchgeführt.\\
Hierbei ist $q_0 = 1-p_0$.\\
Signifikanzniveau $\alpha$\\

\begin{tabular}{llll} 
	$H_0$ 		& $H_1$ 	& Zufallsstreubereich (ZSB) für $k/n$, 	falls $H_0$ zutrifft & $H_0$ verwerfen falls \\
	\toprule
	$p=p_0$ 	& $p \ne p_0$ 	& $[p_0 - z_{1-\alpha/2} \cdot \sqrt{\frac{p_0 q_0}{n}} - \frac{0,5}{n}; \; p_0 + z_{1-\alpha/2} \cdot \sqrt{\frac{p_0 q_0}{n}} + \frac{0,5}{n}]$  & $\frac{k}{n} \not\in$ ZSB \\
	$p \ge p_0$ 	& $p < p_0$ 	& $[p_0 - z_{1-\alpha} \cdot \sqrt{\frac{p_0 q_0}{n}} - \frac{0,5}{n}; \; 1 ]$  & $\frac{k}{n} \not\in$ ZSB\\
	$p \le p_0$ & $p > p_0$ & $[0 ; \; p_0 + z_{1-\alpha} \cdot \sqrt{\frac{p_0 q_0}{n}} + \frac{0,5}{n}]$  & $\frac{k}{n} \not\in$ ZSB
\end{tabular}


\clearpage

\section{Tabellen für Vertrauensbereiche (auch: Konfidenzintervalle)}

\subsection{Vertrauensbereiche für den Erwartungswert $\mu$ einer Normalverteilung bei bekannter Varianz $\sigma^2$}

Die Voraussetzungen sind wie im Abschnitt A. bei den statistischen Tests.\\

Der Vertrauensbereich wird auf Basis einer Stichprobe $x_1, \ldots , x_n$ berechnet, deren arithmetischer Mittelwert $\overline x$ ist.\\

\begin{tabular}{ll} 
Art des Vertrauensbereichs & Vertrauensbereich für $\mu$ \\
	\toprule
zweiseitig & $[\overline x - z_{1-\alpha /2} \cdot \frac{\sigma}{\sqrt{n}} ; \;
\overline x + z_{1-\alpha /2} \cdot \frac{\sigma}{\sqrt{n}}]$ \\
einseitig, nach unten begrenzt & $[\overline x - z_{1-\alpha} \cdot \frac{\sigma}{\sqrt{n}} ; \;
\infty)$ \\
einseitig, nach oben begrenzt & $(-\infty ; \; \overline x + z_{1-\alpha} \cdot \frac{\sigma}{\sqrt{n}}]$ 
\end{tabular}

\subsection{Vertrauensbereiche für den Erwartungswert $\mu$ einer Normalverteilung bei unbekannter Varianz}

Die Voraussetzungen sind wie im Abschnitt B. bei den statistischen Tests.\\

Der Vertrauensbereich wird auf Basis einer Stichprobe $x_1, \ldots , x_n$ berechnet, deren arithmetischer Mittelwert $\overline x$ und deren empirische Standardabweichung $s$ ist.\\

\begin{tabular}{ll}
Art des Vertrauensbereichs & Vertrauensbereich für $\mu$ \\
	\toprule
zweiseitig & $[\overline x - t_{n-1;1-\alpha /2} \cdot \frac{s}{\sqrt{n}} ; \;
\overline x + t_{n-1;1-\alpha /2} \cdot \frac{s}{\sqrt{n}}]$ \\
einseitig, nach unten begrenzt & $[\overline x - t_{n-1;1-\alpha} \cdot \frac{s}{\sqrt{n}} ; \;
\infty)$ \\
einseitig, nach oben begrenzt & $(-\infty ; \;
\overline x + t_{n-1;1-\alpha} \cdot \frac{s}{\sqrt{n}}]$ 
\end{tabular}

\clearpage

\subsection{Vertrauensbereiche für die Differenz $\mu_1 - \mu_2$ der Erwartungswerte zweier Normalverteilungen}

Die Voraussetzungen sind wie im Abschnitt C. bei den statistischen Tests. Der Vertrauensbereich wird auf Basis zweier Stichproben $x_1, \ldots , x_m$ und $y_1, \ldots , y_n$ berechnet.\\

Der arithmetische Mittelwert der ersten Stichprobe ist $\overline x$, der der zweiten ist $\overline y$.

Aus den beiden empirischen Varianzen $s_x^2$ der ersten Stichprobe und $s_y^2$ der zweiten Stichprobe muss zunächst die folgende Hilfsgröße berechnet werden:

\begin{align*}
s_d = \sqrt{(m-1) s_{x}^2 + (n-1) s_{y}^2} \cdot \sqrt{\frac{m+n}{m \cdot n \cdot (m+n-2)}}
\end{align*}



\begin{tabular}{ll}
Art des Vertrauensbereichs & Vertrauensbereich für $\mu_1 - \mu_2$ \\
	\toprule
zweiseitig & $[\overline x - \overline y - t_{m+n-2;1-\alpha /2} \cdot s_d ; \;
\overline x - \overline y + t_{m+n-2;1-\alpha /2} \cdot s_d]$ \\
einseitig, nach unten begrenzt & $[\overline x - \overline y - t_{m+n-2;1-\alpha} \cdot s_d ; \; \infty)$ \\
einseitig, nach oben begrenzt  & $(-\infty ; \;
\overline x - \overline y + t_{m+n-2;1-\alpha} \cdot s_d]$ 
\end{tabular}

\subsection{Vertrauensbereiche für eine Wahrscheinlichkeit $p \;$}

Die Voraussetzungen sind wie im Abschnitt D. bei den statistischen Tests. 

Der Vertrauensbereich wird auf Basis einer Stichprobe vom Umfang $n$ berechnet, bei der das gesuchte Ereignis $k$-mal eingetreten sei, Dann ist $\hat p = \frac{k}{n}$ eine Punktschätzung für die unbekannte Wahrscheinlichkeit. Dieser Schätzwert wird für den Vertrauensbereich benutzt.\\

\begin{tabular}{ll}
Art des Vertrauensbereichs & Vertrauensbereich für $p$ \\
	\toprule
zweiseitig & $[\hat p - z_{1-\alpha/2} \cdot \sqrt{\frac{\hat p \cdot (1- \hat p)}{n}} - \frac{0,5}{n}; \; \hat p + z_{1-\alpha/2} \cdot \sqrt{\frac{\hat p \cdot (1- \hat p)}{n}} + \frac{0,5}{n}]$ \\
einseitig, nach unten begrenzt & $[\hat p - z_{1-\alpha} \cdot \sqrt{\frac{\hat p \cdot (1- \hat p)}{n}} - \frac{0,5}{n}; \; 1 ]$ \\
einseitig, nach oben begrenzt & $[0 ; \; \hat p + z_{1-\alpha} \cdot \sqrt{\frac{\hat p \cdot (1- \hat p)}{n}} + \frac{0,5}{n}]$ 
\end{tabular}


\chapter{Tabellen}
\section{Laplace-Korrespondenzen}
\begin{table}[h]
	\Large
\begin{tabular}{cc|cc}
F(p)&f(t)   & 
F(p)  &f(t) 
\\\toprule
\large{$1  $} &\large{$\delta(t)$}&
$ \frac{1}{(p-a)^2}  $ &\large{$(1 + at) \cdot e^{at}$}
\\\midrule
$\frac{1}{p} $ &\large{$1$}&
$ \frac{p}{(1+ap)^2} $ &\large{$\frac{1}{a^3} \cdot(a-t)\cdot e^{-\frac{t}{a}}$}
\\\midrule
$\frac{1}{p^2} $ &\large{$t$}&
$ \frac{p}{(p-a)(p-b)} $ &$\frac{ae^{at}-be^{bt}}{a-b}$
\\\midrule
$\frac{1}{p^{n+1}} $ &$\frac{t^n}{n!}   $&
$ \frac{p}{(1+ap)(1+bp)} $ &$\frac{ae^{-\frac{t}{b}}-be^{-\frac{t}{a}}}{ab(a-b)}$
\\\midrule
$\frac{1}{p-a} $ &\large{$e^{at}\cdot\sigma{t}$} &
$ \frac{a}{(p-b)^2+a^2} $ &\large{$e^{bt}\cdot\sin(at)$}
\\\midrule
$\frac{1}{1+ap}  $ &$\frac{1}{a}\cdot e^{-\frac{t}{a}} $&
$\frac{p-b}{(p-b)^2+a^2} $ &\large{$e^{bt}\cdot\cos(at)$}
\\\midrule
$\frac{p}{(p^2+a^2)(p^2+b^2)} $ &\large{$ \frac{\cos(bt)-\cos(at)}{a^2-b^2}$}&
$\frac{p^2}{(p^2+a^2)(p^2+b^2)} $ &$\frac{a\cos(at)-b\cos(bt)}{a^2-b^2}$
\\\midrule
$\frac{a}{p^2+a^2} $ &\large{$\sin(at)$}&
$\frac{a}{p^2-a^2}  $ &\large{$\sinh(at)$}
\\\midrule
$\frac{p}{p^2+a^2} $ &\large{$\cos(at)$}&
$\frac{p}{p^2-a^2}  $ &\large{$\cosh(at)$}
\\\midrule
$\frac{1}{p(p-a)}  $ &$\frac{e^{at}-1}{a}  $&
$\frac{p^3}{(p^2+a^2)(p^2+b^2)} $ &$\frac{a^2\cos(at)-b^2\cos(bt)}{a^2-b^2}$
\\\midrule
$\frac{1}{p(1+ap)} $ &\large{$1 - e^{-\frac{t}{a}}$}&
$\frac{a^3}{(p^2+a^2)^2} $ &\large{$\frac{1}{2}\left(\sin(at)-at\cos(at)\right)$}
\\\midrule
$\frac{1}{(p-a)^2} $ &$t \cdot e^{at}   $&
$\frac{ap}{(p^2+a^2)^2} $ &\large{$\frac{t}{2}\cdot\sin(at)$}
\\\midrule
$\frac{1}{(p-a)^{n+1}} $ &$\frac{t^n}{n!} \cdot e^{at} $&
$\frac{ap^2}{(p^2+a^2)^2} $ &\large{$\frac{1}{2}\left(\sin(at)+at\cos(at)\right)$}
\\\midrule
$\frac{1}{(1+pa)^2} $ &$\frac{1}{a^2} \cdot t \cdot e^{-\frac{t}{a}} $&
$\frac{p^3}{(p^2+a^2)^2} $ &\large{$\cos(at) - \frac{at}{2}\cdot\sin(at)$}
\\\midrule
$\frac{1}{(p-a)(p-b)} $ &$\frac{e^{at}-e^{bt}}{a-b}  $&
$\frac{1}{\sqrt{p}}  $ &$\frac{1}{\sqrt{\pi t}}$
\\\midrule
$\frac{1}{(1+ap)(1+bp)} $ &$\frac{e^{-\frac{t}{a}}-e^{-\frac{t}{b}}}{a-b}$&
$\frac{1}{\sqrt{p+a}} $ &$\frac{e^{-at}}{\sqrt{\pi t}}$
\\\midrule
$\frac{ab}{(p^2+a^2)(p^2+b^2)}$ &$\frac{a\sin(bt)-b\sin(at)}{a^2-b^2} $&
$\frac{1}{p^2+2\kappa p+(\kappa^2+\omega^2)} $ &\large{$ \frac{1}{\omega}\cdot e^{-\kappa t} \cdot\sin(\omega t) $}
\end{tabular}
\end{table}


\clearpage

\section{Statistik}
{\bf \slshape \sffamily Quantile $t_{m;q}$ der $t$-Verteilung mit $m$ Freiheitsgraden und}\\
{\bf \slshape \sffamily Quantile $z_{q}$ der Standard-Normalverteilung NV:}
\renewcommand{\arraystretch}{0.9}
\begin{table}[h]
	\small
\begin{tabular}{|c|ccccccc|}
\hline
{} & $q$ & {} & {} & {} & {} & {} & {} \\ \hline
$m$ & 0,8 & 0,9 & 0,95 & 0,975 & 0,99 & 0,995 & 0,999 \\ \hline
1 & 1,376 & 3,078 & 6,314 & 12,706 & 31,821 & 63,656 & 318,289 \\
2 & 1,061 & 1,886 & 2,920 & 4,303 & 6,965 & 9,925 & 22,328 \\
3 & 0,978 & 1,638 & 2,353 & 3,182 & 4,541 & 5,841 & 10,214 \\
4 & 0,941 & 1,533 & 2,132 & 2,776 & 3,747 & 4,604 & 7,173 \\
5 & 0,920 & 1,476 & 2,015 & 2,571 & 3,365 & 4,032 & 5,894 \\ \hline
6 & 0,906 & 1,440 & 1,943 & 2,447 & 3,143 & 3,707 & 5,208 \\
7 & 0,896 & 1,415 & 1,895 & 2,365 & 2,998 & 3,499 & 4,785 \\
8 & 0,889 & 1,397 & 1,860 & 2,306 & 2,896 & 3,355 & 4,501 \\
9 & 0,883 & 1,383 & 1,833 & 2,262 & 2,821 & 3,250 & 4,297 \\
10 & 0,879 & 1,372 & 1,812 & 2,228 & 2,764 & 3,169 & 4,144 \\ \hline
11 & 0,876 & 1,363 & 1,796 & 2,201 & 2,718 & 3,106 & 4,025 \\
12 & 0,873 & 1,356 & 1,782 & 2,197 & 2,681 & 3,055 & 3,930 \\
13 & 0,870 & 1,350 & 1,771 & 2,160 & 2,650 & 3,012 & 3,852 \\
14 & 0,868 & 1,345 & 1,761 & 2,145 & 2,624 & 2,977 & 3,787 \\
15 & 0,866 & 1,341 & 1,753 & 2,131 & 2,602 & 2,947 & 3,733 \\ \hline
16 & 0,865 & 1,337 & 1,746 & 2,120 & 2,583 & 2,921 & 3,686 \\
17 & 0,863 & 1,333 & 1,740 & 2,110 & 2,567 & 2,898 & 3,646 \\
18 & 0,862 & 1,330 & 1,734 & 2,101 & 2,552 & 2,878 & 3,610 \\
19 & 0,861 & 1,328 & 1,729 & 2,093 & 2,539 & 2,861 & 3,579 \\
20 & 0,860 & 1,325 & 1,725 & 2,086 & 2,528 & 2,845 & 3,552 \\ \hline
21 & 0,859 & 1,323 & 1,721 & 2,080 & 2,518 & 2,831 & 3,527 \\
22 & 0,858 & 1,321 & 1,717 & 2,074 & 2,508 & 2,819 & 3,505 \\
23 & 0,858 & 1,319 & 1,714 & 2,069 & 2,500 & 2,807 & 3,485 \\
24 & 0,857 & 1,318 & 1,711 & 2,064 & 2,492 & 2,797 & 3,467 \\
25 & 0,856 & 1,316 & 1,708 & 2,060 & 2,485 & 2,787 & 3,450 \\ \hline
26 & 0,856 & 1,315 & 1,706 & 2,056 & 2,479 & 2,779 & 3,435 \\
27 & 0,855 & 1,314 & 1,703 & 2,052 & 2,473 & 2,771 & 3,421 \\
28 & 0,855 & 1,313 & 1,701 & 2,048 & 2,467 & 2,763 & 3,408 \\
29 & 0,854 & 1,311 & 1,699 & 2,045 & 2,462 & 2,756 & 3,396 \\
30 & 0,854 & 1,310 & 1,697 & 2,042 & 2,457 & 2,750 & 3,385 \\ \hline
35 & 0,852 & 1,306 & 1,690 & 2,030 & 2,438 & 2,724 & 3,340 \\
40 & 0,851 & 1,303 & 1,684 & 2,021 & 2,423 & 2,704 & 3,307 \\
45 & 0,850 & 1,301 & 1,679 & 2,014 & 2,412 & 2,690 & 3,281 \\
50 & 0,849 & 1,299 & 1,676 & 2,009 & 2,403 & 2,678 & 3,261 \\ \hline
60 & 0,858 & 1,296 & 1,671 & 2,000 & 2,390 & 2,660 & 3,232 \\
70 & 0,847 & 1,294 & 1,667 & 1,994 & 2,381 & 2,648 & 3,211 \\
80 & 0,846 & 1,292 & 1,664 & 1,990 & 2,374 & 2,639 & 3,195 \\
90 & 0,846 & 1,291 & 1,662 & 1,987 & 2,368 & 2,632 & 3,183 \\
100 & 0,845 & 1,290 & 1,660 & 1,984 & 2,364 & 2,626 & 3,174 \\ \hline
200 & 0,843 & 1,286 & 1,653 & 1,972 & 2,345 & 2,601 & 3,131 \\
500 & 0,842 & 1,290 & 1,648 & 1,965 & 2,334 & 2,586 & 3,107 \\
{} & {} & {} & {} & {} & {} & {} & {}\\
NV & 0,842 & 1,282 & 1,645 & 1,960 & 2,326 & 2,576 & 3,090 \\ \hline 
\end{tabular}
\end{table}

\noindent
Ablesebeispiele: $t_{20;0,975}=2,086$; $z_{0,995} = 2,576$\\
\noindent
Werte für $q<0,5$ mit den Formeln $t_{m;1-q} = -t_{m;q}$ und $z_{1-q}=-z_q$.\\
Beispiele hierfür: $t_{30;0,1}=-t_{30;0,9}=-1,310$; $z_{0,01}=-z_{0,99}=-2,326$


\clearpage

\noindent
{\bf \slshape \sffamily Verteilungsfunktion $\Phi(z)$ der Standard-Normalverteilung $N(0;1)$:}\\

\begin{table}[h]
	\small
\begin{tabular}{|c|ccccc|ccccc|}
\hline
z & 0 & 1 & 2 & 3 & 4 & 5 & 6 & 7 & 8 & 9 \\ \hline
0,0 & 0,5000 & 0,5040 & 0,5080 & 0,5120 & 0,5160 & 0,5199 & 0,5239 & 0,5279 & 0,5319 & 0,5359 \\
0,1 & 0,5398 & 0,5438 & 0,5478 & 0,5517 & 0,5557 & 0,5596 & 0,5636 & 0,5675 & 0,5714 & 0,5753 \\
0,2 & 0,5793 & 0,5832 & 0,5871 & 0,5910 & 0,5984 & 0,5987 & 0,6026 & 0,6064 & 0,6103 & 0,6141 \\
0,3 & 0,6179 & 0,6217 & 0,6255 & 0,6293 & 0,6331 & 0,6368 & 0,6406 & 0,6443 & 0,6480 & 0,6517 \\
0,4 & 0,6554 & 0,6591 & 0,6628 & 0,6664 & 0,6700 & 0,6736 & 0,6772 & 0,6808 & 0,6844 & 0,6879 \\
0,5 & 0,6915 & 0,6950 & 0,6985 & 0,7019 & 0,7054 & 0,7088 & 0,7123 & 0,7157 & 0,7190 & 0,7224 \\ \hline
0,6 & 0,7257 & 0,7291 & 0,7324 & 0,7357 & 0,7389 & 0,7422 & 0,7454 & 0,7486 & 0,7517 & 0,7549 \\
0,7 & 0,7580 & 0,7611 & 0,7642 & 0,7673 & 0,7704 & 0,7734 & 0,7764 & 0,7794 & 0,7823 & 0,7852 \\
0,8 & 0,7881 & 0,7910 & 0,7939 & 0,7967 & 0,7995 & 0,8023 & 0,8051 & 0,8078 & 0,8106 & 0,8133 \\
0,9 & 0,8159 & 0,8186 & 0,8212 & 0,8238 & 0,8264 & 0,8289 & 0,8315 & 0,8340 & 0,8365 & 0,8389 \\
1,0 & 0,8413 & 0,8438 & 0,8461 & 0,8485 & 0,8508 & 0,8531 & 0,8554 & 0,8577 & 0,8599 & 0,8621 \\ \hline
1,1 & 0,8643 & 0,8665 & 0,8686 & 0,8708 & 0,8729 & 0,8749 & 0,8770 & 0,8790 & 0,8810 & 0,8830 \\
1,2 & 0,8849 & 0,8869 & 0,8888 & 0,8907 & 0,8925 & 0,8944 & 0,8962 & 0,8980 & 0,8997 & 0,9015 \\
1,3 & 0,9032 & 0,9049 & 0,9066 & 0,9082 & 0,9099 & 0,9115 & 0,9131 & 0,9147 & 0,9162 & 0,9177 \\
1,4 & 0,9192 & 0,9207 & 0,9222 & 0,9236 & 0,9251 & 0,9265 & 0,9279 & 0,9292 & 0,9306 & 0,9319 \\
1,5 & 0,9332 & 0,9345 & 0,9357 & 0,9370 & 0,9382 & 0,9394 & 0,9406 & 0,9418 & 0,9429 & 0,9441 \\ \hline
1,6 & 0,9452 & 0,9463 & 0,9474 & 0,9484 & 0,9495 & 0,9505 & 0,9515 & 0,9525 & 0,9535 & 0,9545 \\
1,7 & 0,9554 & 0,9564 & 0,9573 & 0,9582 & 0,9591 & 0,9599 & 0,9608 & 0,9616 & 0,9625 & 0,9633 \\
1,8 & 0,9641 & 0,9649 & 0,9656 & 0,9664 & 0,9671 & 0,9678 & 0,9686 & 0,9693 & 0,9699 & 0,9706 \\
1,9 & 0,9713 & 0,9719 & 0,9726 & 0,9732 & 0,9738 & 0,9744 & 0,9750 & 0,9756 & 0,9761 & 0,9767 \\
2,0 & 0,9772 & 0,9778 & 0,9783 & 0,9788 & 0,9793 & 0,9798 & 0,9803 & 0,9808 & 0,9812 & 0,9817 \\ \hline
2,1 & 0,9821 & 0,9826 & 0,9830 & 0,9834 & 0,9838 & 0,9842 & 0,9846 & 0,9850 & 0,9854 & 0,9857 \\
2,2 & 0,9861 & 0,9864 & 0,9868 & 0,9871 & 0,9875 & 0,9878 & 0,9881 & 0,9884 & 0,9887 & 0,9890 \\
2,3 & 0,9893 & 0,9896 & 0,9898 & 0,9901 & 0,9904 & 0,9906 & 0,9909 & 0,9911 & 0,9913 & 0,9916 \\
2,4 & 0,9918 & 0,9920 & 0,9922 & 0,9925 & 0,9927 & 0,9929 & 0,9931 & 0,9932 & 0,9934 & 0,9936 \\
2,5 & 0,9938 & 0,9940 & 0,9941 & 0,9943 & 0,9945 & 0,9946 & 0,9948 & 0,9949 & 0,9951 & 0,9952 \\ \hline
2,6 & 0,9953 & 0,9955 & 0,9956 & 0,9957 & 0,9959 & 0,9960 & 0,9961 & 0,9962 & 0,9963 & 0,9964 \\
2,7 & 0,9965 & 0,9966 & 0,9967 & 0,9968 & 0,9969 & 0,9970 & 0,9971 & 0,9972 & 0,9973 & 0,9974 \\
2,8 & 0,9974 & 0,9975 & 0,9976 & 0,9977 & 0,9977 & 0,9978 & 0,9979 & 0,9979 & 0,9980 & 0,9981 \\
2,9 & 0,9981 & 0,9982 & 0,9982 & 0,9983 & 0,9984 & 0,9984 & 0,9985 & 0,9985 & 0,9986 & 0,9986 \\
3,0 & 0,9987 & 0,9987 & 0,9987 & 0,9988 & 0,9988 & 0,9989 & 0,9989 & 0,9989 & 0,9990 & 0,9990 \\ \hline
\end{tabular}
\end{table}

\noindent
Ablesebeispiel: $\Phi(0,92) = 0,8212$\\

\noindent
Werte für negatives $z$ mit der Formel $\Phi(-z) = 1-\Phi(z)$, z.B. $\Phi(-1,55) = 1-0,9394 = 0,0606$


%\begin{table}[h]
\begin{tabular}{ll}
	gedämpfte Schwingung & Eine Schwingung ist gedämpft, wenn sie mit $\e^{-x}$ multipliziert wird\\
	$\omega$ & Elementarereignis\\
	$A$  & Ausgang, (günstiges) Ereignis, der gesuchte Wert\\
	$N$  & Anzahl von Möglichkeiten\\
	$|A|$& Anzahl günstiger Fälle / Häufigkeit\\
	$|\Omega|$ & Anzahl möglicher Fälle / Summe der Häufigkeiten\\
	$P(A)=\frac{|A|}{|\Omega|}$ & Wahrscheinlichkeit Für Ausgang / Ereignis A\\
	$f(x_i)$ & Wahrscheinlichkeit / Diskrete Dichte / $P(X=x)$\\
	$F(x)$ & Verteilungsfunktion, aufsummierte W.-Dichte\\
	$\mu$ & arithmetisches Mittel d. W\\
	$\sigma^2=Var(X)$ & Varianz\\
	$\sigma=\sqrt{\sigma^2}$ & Standardabweichung\\
	$X(\omega)$ & Zufallsvariable, Anzahl d. Treffer pro $\omega$\\
	& (wie oft kommt Zahl in 5 Münzwurfen)\\
	$\Omega$ & Anz. d. Durchgänge für $\omega$\\
	& (wie oft werden 5 Münzen geworfen?) \\
	$\alpha$ & Signifikanzniveau (Unsicherheit)\\
	$\Phi$ & Standardnormalverteilung\\
	Regressionsgerade & Beträge der Abstände der Messwerte von der Gerade\\
	Zufallsexperiment & Ausgang offen und zufällig\\
	$\binom{n}{k}$ & Binomialkoeffizient\\
	BV & Binomialverteilung\\
	$\hat{p}=\frac{k}{n}$ & Punktschätzer für die Wahrscheinlichkeit\\
	$\notin$ & nicht Element von \\ 	
	$s^2$ & empirische Standardabweichung\\
	$s$ & empirische Varianz\\
	$\mu$ & Erwartungswert / Sollwert\\
	$\Phi$ & Standardnormalverteilung\\
	$\Phi\left( \frac{\alpha-\mu}{\sigma} \right)$ &  \\
	$f(x)$  & Wahrscheinlichkeitsdichtefunktion\\
	$F(x)$ & Verteilungsfunktion\\
	$x_{0,95}$ & 95\%-Quantil (der Normalverteilung)\\
	$P(X < x_{0,95}) = F(x_{0,95}) = 0,95$ &\\
	ZSB & Zufallsstreubereich\\
\end{tabular}
\end{table}



% Hier kommmen Tabellen und Bilder
%\listoftables
%\listoffigures

\end{document}
%% Dokument ENDE %%%%%%%%%%%%%%%%%%%%%%%%%%%%%%%%%%%%%%%%%%%%%%%%%%%%%%%%%%

